% concatenated from main.tex, theorems.tex, definitions.tex
%%%%%%%% ICML 2026 EXAMPLE LATEX SUBMISSION FILE %%%%%%%%%%%%%%%%%

\newif\ifarxiv
\arxivtrue  % Preprint publication. Comment for Journal version

\ifarxiv
  \documentclass[12pt]{article}
\else
  \documentclass{article}
\fi

% \usepackage{etoolbox}
% \newtoggle{preprint}
%\togglefalse{preprint}
% \toggletrue{preprint}
\usepackage{csvsimple}
\usepackage{natbib}
\ifarxiv
\setcitestyle{authoryear,round,citesep={;},aysep={,},yysep={;}}
\renewcommand{\cite}[1]{\citep{#1}}
\fi
\ifarxiv
\usepackage[T1]{fontenc}   % switch from OT1 to T1 encoding
\usepackage[utf8]{inputenc}% if your source is UTF-8 (likely)
\usepackage{textcomp}      % provides \guilsinglleft \guilsinglright
\usepackage{lmodern}       % Latin Modern fonts with T1 glyphs (optional but good)
% \usepackage[colorlinks=true, citecolor=cyan, linkcolor=cyan, urlcolor=cyan]{hyperref}

% Prevent words to overflow over the margin
% \sloppy  % Alternative command
\emergencystretch 3em

\usepackage{fullpage}
% \setlength{\parindent}{0pt}

% Images and plots
\usepackage{graphicx}
\usepackage{subcaption}
\usepackage{pifont}% http://ctan.org/pkg/pifont
\usepackage{pgfplots}

% Math
\usepackage{amsmath,amssymb,amsthm,enumitem,mathtools}
% Theorem-like environments (share numbering)
\newtheorem{theorem}{Theorem}[section]     % Theorem 2.1, 2.2, ...
\newtheorem{lemma}[theorem]{Lemma}         % Lemma 2.3, ...
\newtheorem{corollary}[theorem]{Corollary} % Corollary 2.4, ...
\newtheorem{proposition}[theorem]{Proposition}
\newtheorem{assumption}[theorem]{Assumption}

% Definition-style
\theoremstyle{definition}
\newtheorem{definition}[theorem]{Definition}
\newtheorem{problem}[theorem]{Problem}
\newtheorem{example}[theorem]{Example}
\newtheorem{exercise}[theorem]{Exercise}

% Remark-style
\theoremstyle{remark}
\newtheorem{remark}[theorem]{Remark}

\renewcommand{\qedsymbol}{$\blacksquare$}
\renewcommand\arraystretch{1.2}
\usepackage{multirow}
\usepackage{algorithm}% http://ctan.org/pkg/algorithms
\usepackage[noend]{algpseudocode} % Avoid "end" and make it look cleaner

% Tables
\usepackage{csvsimple}	% reading CSV files in tables
\usepackage{booktabs}   % Nicer Tables
\usepackage{adjustbox}  % To adjust table length

% Hyperreferences
\usepackage[colorlinks=true, citecolor=cyan, linkcolor=cyan, urlcolor=cyan]{hyperref}

% Acronyms
\usepackage[xindy, acronyms]{glossaries}   % Glossaries and Acronyms
\makeglossaries

% Todos
\newcommand{\snote}[1]{}
\newcommand{\bnote}[1]{}
% Comment to hide notes
\renewcommand{\snote}[1]{\textcolor{blue}{\textbf{\small [S: #1]}}}
\renewcommand{\bnote}[1]{\textcolor{blue}{\textbf{\small [B: #1]}}}

% Definitions

\newcommand{\paramdist}{\mathbf{Q}}
\newcommand{\paramspace}{\Theta}
\newcommand{\param}{\theta}
% \newcommand{\counter}{\mathcal{N}}  % removed: inlined as \#\{...\}

\newcommand{\reals}{\mathbf{R}}
\newcommand{\cardinality}{\#}

\newcommand{\st}{\tau}
\newcommand{\Z}{\mathbf{Z}}
\newcommand{\Prob}{\mathbf{P}}
\newcommand{\MIP}{\operatorname{MIP}}

\newcommand{\UB}{U}
\newcommand{\LB}{L}

\newcommand{\X}{\hat{Y}}
\newcommand{\x}{\hat{y}}
\newcommand{\ov}{z^\star}
\newcommand{\covars}{X}

\newcommand{\nn}{h}
\newcommand{\activ}{\phi}

\newcommand{\loss}{\mathcal{L}}
\newcommand{\E}{\mathbf{E}}
\newcommand{\dataset}{\mathcal{D}}

\newcommand{\varspace}{\mathcal{X}}
\newcommand{\vars}{x}
\newcommand{\transpose}{T}

\newcommand{\algo}{\mathcal{A}}

\newcommand{\truegap}{g}
\newcommand{\gap}{g^{\mathrm{alg}}}
\newcommand{\predgap}{\hat{g}}

\newcommand{\tol}{\epsilon}
\newcommand{\errorprob}{\alpha}
\newcommand{\pst}{\hat{\tau}}
\newcommand{\approxerrorprob}{\hat{k}}

\newcommand{\image}[1]{\mathcal{I}(#1)}

\newcommand{\calibration}{\mathcal{C}}

\newcommand{\ect}{\emph{ect}}
\newcommand{\ie}{\emph{i.e.}}
\newcommand{\eg}{\emph{e.g.}}
\newcommand{\termov}{s}
\newcommand{\lt}{\kappa}
\newcommand{\zbd}{{S_{\text{max}}}}
\newcommand{\classif}{\rho}

\newcommand{\rademacher}{\mathcal{R}}
\newcommand{\vc}{\nu}
\newcommand{\indicator}{\mathbf{1}}


\newcommand{\test}{\mathcal{T}}

% EXPERIMENTS
\newcommand{\LSTMLAYERS}{2}
\newcommand{\LSTMHIDDEN}{200}
\newcommand{\FFLAYERS}{2}
\newcommand{\FFHIDDEN}{200}
\newcommand{\NTRAIN}{9000}
\newcommand{\NTEST}{100}
\newcommand{\NCALIB}{1000}

\newcommand{\faclocx}{x}
\newcommand{\faclocy}{y}
\newcommand{\faclocz}{z}
\newcommand{\weight}{w}
\newcommand{\ncalib}{c}
\newcommand{\rollmin}{\bar{g}}

\newcommand{\train}{\mathcal{D}}
\newcommand{\ntrain}{d}
\newcommand{\ntest}{l}
\newcommand{\runtime}{T}
% \newcommand{\truerelgap}{r}  % removed: using only relative gap via \truegap
\newcommand{\tbd}{T_{\max}}

% Camera-ready: render \change content in plain black. Keep the wrappers in
% source so we still have a diff record of revision-tracked passages.
\newcommand{\change}[1]{#1}
% \newcommand{\change}[1]{{\color{red}#1}}  % uncomment to highlight revisions

% \title{Learning Sequential Mixed Integer Programming Policies}
\title{Conformal Prediction for Early Stopping in Mixed Integer Optimization}
\author{Stefan Clarke and Bartolomeo Stellato}
\date{\small Department of Operations Research and Financial Engineering\\[.5em] Princeton University\\[.5em]\today}
\else
% Todos
\newcommand{\snote}[1]{}
\newcommand{\bnote}[1]{}

% Comment to hide notes
\renewcommand{\snote}[1]{\textcolor{blue}{\textbf{\small [Stefan: #1]}}}
\renewcommand{\bnote}[1]{\textcolor{blue}{\textbf{\small [B: #1]}}}
% Definitions

\newcommand{\paramdist}{\mathbf{Q}}
\newcommand{\paramspace}{\Theta}
\newcommand{\param}{\theta}
% \newcommand{\counter}{\mathcal{N}}  % removed: inlined as \#\{...\}

\newcommand{\reals}{\mathbf{R}}
\newcommand{\cardinality}{\#}

\newcommand{\st}{\tau}
\newcommand{\Z}{\mathbf{Z}}
\newcommand{\Prob}{\mathbf{P}}
\newcommand{\MIP}{\operatorname{MIP}}

\newcommand{\UB}{U}
\newcommand{\LB}{L}

\newcommand{\X}{\hat{Y}}
\newcommand{\x}{\hat{y}}
\newcommand{\ov}{z^\star}
\newcommand{\covars}{X}

\newcommand{\nn}{h}
\newcommand{\activ}{\phi}

\newcommand{\loss}{\mathcal{L}}
\newcommand{\E}{\mathbf{E}}
\newcommand{\dataset}{\mathcal{D}}

\newcommand{\varspace}{\mathcal{X}}
\newcommand{\vars}{x}
\newcommand{\transpose}{T}

\newcommand{\algo}{\mathcal{A}}

\newcommand{\truegap}{g}
\newcommand{\gap}{g^{\mathrm{alg}}}
\newcommand{\predgap}{\hat{g}}

\newcommand{\tol}{\epsilon}
\newcommand{\errorprob}{\alpha}
\newcommand{\pst}{\hat{\tau}}
\newcommand{\approxerrorprob}{\hat{k}}

\newcommand{\image}[1]{\mathcal{I}(#1)}

\newcommand{\calibration}{\mathcal{C}}

\newcommand{\ect}{\emph{ect}}
\newcommand{\ie}{\emph{i.e.}}
\newcommand{\eg}{\emph{e.g.}}
\newcommand{\termov}{s}
\newcommand{\lt}{\kappa}
\newcommand{\zbd}{{S_{\text{max}}}}
\newcommand{\classif}{\rho}

\newcommand{\rademacher}{\mathcal{R}}
\newcommand{\vc}{\nu}
\newcommand{\indicator}{\mathbf{1}}


\newcommand{\test}{\mathcal{T}}

% EXPERIMENTS
\newcommand{\LSTMLAYERS}{2}
\newcommand{\LSTMHIDDEN}{200}
\newcommand{\FFLAYERS}{2}
\newcommand{\FFHIDDEN}{200}
\newcommand{\NTRAIN}{9000}
\newcommand{\NTEST}{100}
\newcommand{\NCALIB}{1000}

\newcommand{\faclocx}{x}
\newcommand{\faclocy}{y}
\newcommand{\faclocz}{z}
\newcommand{\weight}{w}
\newcommand{\ncalib}{c}
\newcommand{\rollmin}{\bar{g}}

\newcommand{\train}{\mathcal{D}}
\newcommand{\ntrain}{d}
\newcommand{\ntest}{l}
\newcommand{\runtime}{T}
% \newcommand{\truerelgap}{r}  % removed: using only relative gap via \truegap
\newcommand{\tbd}{T_{\max}}

% Camera-ready: render \change content in plain black. Keep the wrappers in
% source so we still have a diff record of revision-tracked passages.
\newcommand{\change}[1]{#1}
% \newcommand{\change}[1]{{\color{red}#1}}  % uncomment to highlight revisions

% Recommended, but optional, packages for figures and better typesetting:
\usepackage{microtype}
\usepackage{graphicx}
\usepackage{subcaption}
\usepackage{booktabs} % for professional tables
\usepackage{multirow}

% hyperref makes hyperlinks in the resulting PDF.
% If your build breaks (sometimes temporarily if a hyperlink spans a page)
% please comment out the following usepackage line and replace
% \usepackage{icml2026} with \usepackage[nohyperref]{icml2026} above.
\usepackage{hyperref}

% Attempt to make hyperref and algorithmic work together better:
\newcommand{\theHalgorithm}{\arabic{algorithm}}

% Use the following line for the initial blind version submitted for review:
% \usepackage{icml2026}
% \usepackage[accepted]{icml2026}

% For preprint, use
% \usepackage[preprint]{icml2026}

% If accepted, instead use the following line for the camera-ready submission:
\usepackage[accepted]{icml2026}
\def\isaccepted{}

\usepackage{amsmath}
\usepackage{amssymb}
\usepackage{mathtools}
\usepackage{amsthm}

% if you use cleveref..
\usepackage[capitalize,noabbrev]{cleveref}

% Balance columns on final page
\usepackage{balance}

%%%%%%%%%%%%%%%%%%%%%%%%%%%%%%%%
% THEOREMS
%%%%%%%%%%%%%%%%%%%%%%%%%%%%%%%%
\theoremstyle{plain}
\newtheorem{theorem}{Theorem}[section]
\newtheorem{proposition}[theorem]{Proposition}
\newtheorem{lemma}[theorem]{Lemma}
\newtheorem{corollary}[theorem]{Corollary}
\theoremstyle{definition}
\newtheorem{definition}[theorem]{Definition}
\newtheorem{assumption}[theorem]{Assumption}
\theoremstyle{remark}
\newtheorem{remark}[theorem]{Remark}

% Todonotes is useful during development; simply uncomment the next line
%    and comment out the line below the next line to turn off comments
%\usepackage[disable,textsize=tiny]{todonotes}
\usepackage[textsize=tiny]{todonotes}

% The \icmltitle you define below is probably too long as a header.
% Therefore, a short form for the running title is supplied here:
\icmltitlerunning{Conformal Prediction for Early Stopping in Mixed Integer Optimization}
\fi

\begin{document}

\ifarxiv
\maketitle
\else
\twocolumn[
  \icmltitle{Conformal Prediction for Early Stopping in Mixed Integer Optimization}

  % It is OKAY to include author information, even for blind submissions: the
  % style file will automatically remove it for you unless you've provided
  % the [accepted] option to the icml2026 package.

  % List of affiliations: The first argument should be a (short) identifier you
  % will use later to specify author affiliations Academic affiliations
  % should list Department, University, City, Region, Country Industry
  % affiliations should list Company, City, Region, Country

  % You can specify symbols, otherwise they are numbered in order. Ideally, you
  % should not use this facility. Affiliations will be numbered in order of
  % appearance and this is the preferred way.
  \icmlsetsymbol{equal}{*}

  \begin{icmlauthorlist}
    \icmlauthor{Stefan Clarke}{princeton}
    \icmlauthor{Bartolomeo Stellato}{princeton}

  \end{icmlauthorlist}

  \icmlaffiliation{princeton}{Operations Research and Financial Engineering, Princeton University, Princeton, USA}

  \icmlcorrespondingauthor{Bartolomeo Stellato}{bstellato@princeton.edu}

  % You may provide any keywords that you find helpful for describing your
  % paper; these are used to populate the "keywords" metadata in the PDF but
  % will not be shown in the document
  \icmlkeywords{Conformal Prediction, Mixed-Integer Optimization, Learning for Optimization}

  \vskip 0.3in
]

% this must go after the closing bracket ] following \twocolumn[ ...

% This command actually creates the footnote in the first column listing the
% affiliations and the copyright notice. The command takes one argument, which
% is text to display at the start of the footnote. The \icmlEqualContribution
% command is standard text for equal contribution. Remove it (just {}) if you
% do not need this facility.

% Use ONE of the following lines. DO NOT remove the command.
% If you have no special notice, KEEP empty braces:
\printAffiliationsAndNotice{}  % no special notice (required even if empty)
% Or, if applicable, use the standard equal contribution text:
% \printAffiliationsAndNotice{\icmlEqualContribution}
\fi


\begin{abstract}
Mixed-integer optimization solvers often find optimal solutions early in the search, yet spend the majority of computation time proving optimality.
We exploit this by learning when to terminate solvers early on distributions of similar problem instances.
Our method trains a neural network to estimate the true optimality gap from the solver state, then uses conformal prediction to calibrate a stopping threshold with rigorous probabilistic guarantees on solution quality.
On \change{six} problem families from the Distributional MIPLIB library, our method reduces solve time by over 60\% while guaranteeing $0.1\%$-optimal solutions with 95\% probability \change{for new instances drawn from the same distribution}.
\end{abstract}

\section{Introduction}
Mixed-integer optimization problems arise in robotics~\cite{Shin2022DesignOK, Halsted2021ASO}, transportation~\cite{Toth2002TheVR, Bertsimas2019OnlineVR}, power systems~\cite{VanHentenryck2021MachineLF, Sarkar2018ReactivePM}, and control~\cite{Marcucci2019WarmSO}.
Many applications require decisions in real time, yet mixed-integer programming is NP-hard~\cite{garey1979computers}, and even modern solvers~\cite{gurobi, Qi2015ParallelizingTD} may be too slow for time-critical settings.

The standard method for solving mixed-integer programs is branch-and-bound~\cite{Land1960AnAM}, often combined with cutting planes to form branch-and-cut~\cite{Mitchell2002BranchandCutAF}.
As the algorithm progresses it maintains upper and lower bounds on the optimal objective value; when these bounds meet, the solver terminates with a provably optimal solution.
Heuristics help find good feasible solutions early~\cite{Fischetti2011HeuristicsIM, Angioni2025NeuralCO}, quickly reducing the upper bound, but closing the lower bound to certify optimality typically requires many more iterations.
The result is that the solver often discovers an optimal solution long before it can prove optimality (see Figure~\ref{fig:bbsolve}).
\begin{figure}[t]
    \centering
    \ifarxiv
    \includegraphics[width=0.55\linewidth, trim={0cm 0.2cm 0 0.2cm}]{bounds_plot.pdf}
    \else
    \includegraphics[width=1.0\linewidth, trim={0cm 0.2cm 0 0.2cm}]{bounds_plot.pdf}
    \fi
    \caption{Bounds for one solve of one instance of the Distributional MIPLIB optimal transmission switching problem family~\cite{Huang2024DistributionalMA}.
    The optimal solution is found after about 57 seconds, but the solve does not terminate until nearly 120 seconds have passed because the gap is not within the tolerance.
    Our method with a learned lower-bound terminates the solve much earlier, around 57 seconds.
    The upper-bound remains large until around 20 seconds, after which it quickly decreases as the solver finds feasible solutions.
    % We add guarantees that the solution returned by the learned solver is within a specified tolerance of the true optimal solution with high probability in this paper.
    }
    \label{fig:bbsolve}
\end{figure}

Machine learning has emerged as a promising tool to accelerate mixed-integer optimization.
\citet{Bengio2021MachineLF} survey the landscape, while foundational work includes learning combinatorial optimization algorithms over graphs~\cite{Dai2017LearningCO} and learning branching policies~\cite{Khalil2016LearningTB}.
Subsequent efforts have targeted branching decisions~\cite{Scavuzzo2022LearningTB}, primal heuristics~\cite{Shen2021LearningPH, Angioni2025NeuralCO}, solver parameters~\cite{Patel2023ProgressivelySA}, and cutting planes~\cite{Wang2023LearningCS, Dragotto2023DifferentiableCL}.
These methods aim to reach a provably optimal solution faster by improving decisions within the solver.

In this paper we take a different approach: rather than speeding up the search for optimality, we learn when to stop early.
We train a neural network to predict the true optimality gap from the solver state and use conformal prediction~\cite{Angelopoulos2021AGI} to calibrate a stopping threshold with rigorous probabilistic guarantees on solution quality.
On benchmarks from the Distributional MIPLIB library~\cite{Huang2024DistributionalMA}, our method achieves significant speedups while provably returning near-optimal solutions with high probability.

\section{Related works}

\subsection{Learning to accelerate mixed-integer optimization}
Machine learning has emerged as a powerful tool to accelerate mixed-integer optimization~\cite{Bengio2021MachineLF, Zhang2022ASF}.
The most studied application is learning branching policies.
\citet{Khalil2016LearningTB} pioneered imitation learning to approximate strong branching.
\citet{Gasse2019ExactCO} introduced graph neural networks that exploit the bipartite variable-constraint structure of MILPs, inspiring subsequent work on hybrid models combining graph neural networks and multi-layer perceptrons~\cite{Gupta2020HybridMF} and tree-based formulations~\cite{Scavuzzo2022LearningTB}.
\citet{Scavuzzo2024MLaugmentedBB} provide a comprehensive survey of machine learning for branch-and-bound.
Beyond branching, researchers have applied learning to cutting plane selection~\cite{Wang2023LearningCS, Huang2021LearningTS, Dragotto2023DifferentiableCL}, solver configuration~\change{\cite{Patel2023ProgressivelySA, Gao2024NeuralSS, Liu2024L2PMIP}}, and primal heuristics~\change{\cite{Shen2021LearningPH, Nair2020SolvingMI, Angioni2025NeuralCO, Han2023GNNPredict}}.
Some methods learn to predict solutions directly without solving~\cite{Sun2023DIFUSCOGD}.
\change{Apollo-MILP~\cite{Liu2025ApolloMILP} uses an alternating prediction-correction loop to fix high-confidence variable values and reduce the problem dimension, using an uncertainty bound to identify reliable predictions; however, it targets primal solution quality rather than certifying when to stop the solver.}

All these approaches aim to reach a provably optimal solution faster by improving decisions within the solver.
Our method takes a fundamentally different approach: we observe that branch-and-bound often finds an optimal solution long before it can prove optimality (Figure~\ref{fig:bbsolve}), and we learn when to stop early rather than speeding up the proof.

\subsection{Conformal prediction}
Conformal prediction provides distribution-free uncertainty quantification with finite-sample guarantees~\cite{Angelopoulos2021AGI}.
It has been applied across domains including natural language processing~\cite{Campos2024ConformalPF}, time-series forecasting~\cite{Stankeviit2021ConformalTF}, medical imaging~\cite{Vazquez2022ConformalPI, Fayyad2023EmpiricalVO}, and autonomous driving~\cite{Doula2024ConformalPF}.
Recent work has extended conformal prediction to optimization under uncertainty, providing calibration guarantees for robust optimization~\cite{Jia2024EndtoEndCC}.
\citet{Clarke2025LearningBasedHA} used conformal prediction to bound the suboptimality of learned heuristics for parametric MIPs.
\change{This work differs from~\cite{Clarke2025LearningBasedHA} in that we learn a stopping criterion for the exact solver rather than validating an independent heuristic. In the previous work, the learned heuristic and the solution-finding heuristic are independent of one another, and the use of conformal prediction is to give bounds on the error of the learned heuristic. In this work, we explicitly set an acceptable error probability and magnitude of error, and use conformal prediction to calibrate a stopping criterion that achieves these guarantees.}

This work is arguably the first to use conformal prediction to accelerate exact MIP solving.
Rather than validating an independent heuristic, in this work we learn a termination criterion that stops the solver early while providing probabilistic guarantees on solution quality. 

\begin{figure}
    \centering
    \ifarxiv
    \includegraphics[width=0.55\linewidth, trim={0cm 12cm 27cm 0cm}]{method_diagram.pdf}
    \else
    \includegraphics[width=\linewidth, trim={0cm 12cm 27cm 0cm}]{method_diagram.pdf}
    \fi
    \caption{Diagram of the proposed conformal prediction method for accelerating MIP solvers by learning to terminate early.
    The LSTM predicts the suboptimality-gap and the solver terminates and returns the current best solution when the predicted gap is below a threshold. The threshold is chosen by conformal prediction.}\label{fig:method_overview}
\end{figure}


\section{Parametric mixed integer optimization}
This section formalizes the problem setting.
We define parametric mixed-integer programs, describe how solvers maintain bounds over time, and show that standard termination criteria are conservative: the solver may hold a near-optimal solution long before it can prove optimality.

\subsection{Parametric mixed-integer programs} \label{sec:parametric}
Let $\paramdist$ be a distribution supported on a set $\paramspace \subseteq \reals^p$.
For each $\param \in \paramspace$ let $\MIP_\param$ be the optimization problem,
\begin{equation} \label{eq:master} \tag{$\MIP_\param$}
    \ov(\param) = \begin{array}[t]{ll}
         \text{minimize}  & {c(\param)}^\transpose \vars \\
         \text{subject to} & \vars \in \varspace(\param),
    \end{array}
\end{equation}
where $\varspace(\param) \subseteq \reals^m$ is the feasible set, which may include integrality constraints on some variables, and $c(\param) \in \reals^m$ is the cost vector.
We consider repeated solves of~\eqref{eq:master} for independent draws of $\param \sim \paramdist$.
We write $\vars^\star(\param)$ for an optimal solution to~\eqref{eq:master}.

\subsection{Solver behavior over time} \label{sec:solver-time}
Consider a fixed $\param \in \paramspace$.
An algorithm~$\algo(\param)$ designed to solve~\eqref{eq:master} can be viewed as a function of time.
At time $t \geq 0$ the algorithm holds an upper bound $\UB_\param(t)$ and a lower bound $\LB_\param(t)$ such that $\LB_\param(t) \le \ov(\param) \le \UB_\param(t)$.
If $\UB_\param(t) < \infty$ then it also holds a feasible solution $\vars_\param(t) \in \varspace(\param)$ such that $c(\param)^\transpose \vars_\param(t) = \UB_\param(t).$
We assume that $\param$ and $t$ uniquely identify the algorithm state; this approximation holds when problems are solved on the same hardware.

For the solver to eventually certify optimality, the bounds must converge.
We formalize this as the following assumption, which holds for any complete branch-and-bound algorithm.
\begin{assumption} \label{ass:terminal}
    $\UB_\param(t) - \LB_\param(t) \rightarrow 0$ as $t \rightarrow \infty$.
\end{assumption}
Define the~\emph{true optimality gap} at time $t > 0$ for algorithm $\algo(\param)$ by
\begin{equation*}
    \truegap_\param(t) = \frac{\UB_\param(t) - \ov(\param)}{|\ov(\param)|},
\end{equation*}
assuming $\ov(\param) \neq 0$.
Assumption~\ref{ass:terminal} implies that $\truegap_\param(t) \rightarrow 0$ as $t \rightarrow \infty$.
The true optimality gap is unknown during the solve but can be computed afterward since the solver returns~$\ov(\param)$.

\subsection{Standard termination criteria} \label{sec:deterministic}
Solvers typically terminate when the optimality gap
\begin{equation*}
    \gap_\param(t) = \frac{\UB_\param(t) - \LB_\param(t)}{|\LB_\param(t)|},
\end{equation*}
drops below a threshold $\tol > 0$.
We define the $\tol$-optimality stopping time for algorithm $\algo$ as
\begin{equation*}
    \st_\tol(\param) = \inf\{t \ge 0 \mid \gap_\param(t) \le \tol\}.
\end{equation*}
At time $\st_\tol(\param)$, the best feasible solution $\vars_\param(\st_\tol(\param))$ is $\tol$-optimal because $\truegap_\param(\st_\tol(\param)) \le \gap_\param(\st_\tol(\param)) \le \tol$, where the first inequality uses $\LB_\param(t) \le \ov(\param)$.
However, the first inequality may not be tight.
At time $\st_\tol(\param)$, the algorithm may have held the solution $\vars_\param(\st_\tol(\param))$ for some time already.
The algorithm could have terminated earlier and still returned an $\tol$-optimal solution.
This observation is the foundation of our approach.

\subsection{Probabilistic termination criteria} \label{sec:probabilistic}
At the deterministic stopping time $\st_\tol(\param)$, we have
\begin{equation*}
\Prob[\truegap_\param(\st_\tol(\param)) \le \tol] = 1,
\end{equation*}
where the probability is over $\param \sim \paramdist$.
We relax this by allowing a small error probability $\errorprob \in (0,1)$, seeking stopping times $\pst(\param)$ satisfying
\begin{align} \label{eq:prob}
\Prob[\truegap_\param(\pst(\param)) \le \tol] \ge 1 - \errorprob.
\end{align}
Allowing error probability $\errorprob$ enables much earlier termination than deterministic criteria.
Section~\ref{sec:learning} describes how we construct such stopping times by learning to approximate $\truegap_\param(t)$.

\section{Learning termination criteria} \label{sec:learning}
This section develops a learned stopping criterion with rigorous probabilistic guarantees.
The key idea is to train a predictor $\predgap_\param(t)$ that approximates the true optimality gap $\truegap_\param(t)$, then use conformal prediction to calibrate a threshold $\lt$ such that terminating when $\predgap_\param(t) \le \lt$ satisfies~\eqref{eq:prob}.

Given predictor $\predgap_\param(t)$, define the stopping time
\begin{equation*}
    \pst_\lt(\param) = \inf \{t \ge 0 \mid \predgap_\param(t) \le \lt\},
\end{equation*}
where $\lt$ depends on the desired tolerance $\tol$ and error probability $\errorprob$.
We assume $\predgap_\param(t) \to 0$ as $t \to \infty$ for all $\param \in \paramspace$; Section~\ref{sec:training} describes how the predictor enforces this property.
\subsection{Conformal prediction for termination} \label{sec:conformal}
We develop a conformal prediction framework to choose $\lt$ given $\tol$ and $\errorprob$ such that~\eqref{eq:prob} holds.
The construction requires two technical definitions.

Define the \emph{left-inverse} of a function $f:\reals_+ \rightarrow \reals_+$ by
\begin{equation*}
    f^{-1}(x) = \begin{cases}
        \inf\{t \ge 0 \mid f(t) \le x\} & x \le \sup f \\
        \infty & \text{otherwise}.
    \end{cases}
\end{equation*}
Define the \emph{rolling minimum} of the predictor by
\begin{equation*}
    \rollmin_\param(t) = \min_{0 \le s \le t} \predgap_\param(s).
\end{equation*}
The function $\rollmin_\param(t)$ is nonincreasing in $t$.
Let $\predgap_\param^{-1}$ and $\truegap_\param^{-1}$ denote the left-inverses of $\predgap_\param$ and $\truegap_\param$.
For $k \ge 0$ define $\pst_k(\param) = \predgap^{-1}_\param(k)$.
The following theorem, based on conformal prediction theory~\cite{Angelopoulos2021AGI}, shows how to choose $\lt$.
\begin{theorem} \label{thm:conformal}
    Let $\ncalib \in \Z_+$ and $n \in \{1, \dots, \ncalib\}$.
    Let $\param_1, \dots, \param_\ncalib, \param_{\ncalib+1} \sim \paramdist$ be independent.
    Let
    \begin{equation*}
        \lt = \sup \{ k \ge 0 \mid \#\{i \in \{1, \dots, \ncalib\} \mid \truegap_{\param_i}(\pst_k(\param_i)) \le \tol\} \ge n\}.
    \end{equation*}
    Then $\Prob[\truegap_{\param_{\ncalib+1}}(\pst_{\lt}(\param_{\ncalib+1})) \le \tol] \ge \change{n}/(\ncalib + 1)$.
\end{theorem}
The proof relies on the following lemma.
\begin{lemma} \label{lemma:ordering}
    Let $Z_1, \dots, Z_{\ncalib + 1}$ be iid random variables in $\reals$.
    Let $Z_{[1]}, \dots, Z_{[\ncalib]}$ denote $Z_1, \dots, Z_\ncalib$ arranged in increasing order.
    For any $n \in \{1, \dots, \ncalib\}$,
    $$
    \Prob[Z_{\ncalib + 1} \ge Z_{[\ncalib + 1 - n]}] \ge n / (\ncalib + 1).
    $$
\end{lemma}
\begin{proof}
    Let $m$ be the rank of $Z_{\ncalib+1}$ among $Z_1, \dots, Z_{\ncalib + 1}$ with random tie-breaking.
    Since $m \ge \ncalib + 2 - n$ implies $Z_{\ncalib + 1} \ge Z_{[\ncalib + 1 - n]}$,
    \begin{align*}
        \Prob[Z_{\ncalib + 1} \ge Z_{[\ncalib + 1 - n]}] &\ge \Prob[m \ge \ncalib + 2 - n] \\
        & = \sum_{j=\ncalib + 2 - n}^{\ncalib + 1} \Prob[m = j], \\
        &= \sum_{j=\ncalib + 2 - n}^{\ncalib + 1} \frac{1}{\ncalib + 1} \\
        &= \frac{n}{\ncalib + 1}.
    \end{align*}
    where the second-last equality holds by exchangeability of $Z_1, \dots, Z_{\ncalib+1}$.
\end{proof}
% \begin{lemma} \label{lemma:rollmin}
%     For any $\param$, $t$ we have $\rollmin_\param(t) \le \tol \iff \truegap^{-1}(\tol) \le t$.
% \end{lemma}
% \begin{proof}
% \end{proof}
\begin{proof}[Proof of Theorem~\ref{thm:conformal}]
    By definition, $\lt$ is the largest $k \ge 0$ such that at least $n$ of $\{\truegap_{\param_i}(\pst_k(\param_i))\}_{i=1}^\ncalib$ are at most $\tol$.
    Since $\truegap_\param(t)$ is nonincreasing,
    $$
    \truegap_\param(t) \le \tol \iff \truegap_\param^{-1}(\tol) \le t.
    $$
    We also have $\rollmin_\param^{-1} = \predgap_\param^{-1}$, so
    \begin{align}
        \rollmin_\param(t) \ge k \iff \rollmin_\param^{-1}(k) \ge t \iff \predgap_\param^{-1}(k) \ge t. \label{eq:confproof1}
    \end{align}
    % Let 
    % $$
    % m = \#\{i \in \{1, \dots, \ncalib\} \mid \truegap_{\param_i}(\pst_k(\param_i)) \le \tol\}.
    % $$
    Observe that,
    \begin{align*}
    \truegap_{\param_i}^{-1}(\tol) \le \pst_k(\param_i) &\iff \truegap_{\param_i}^{-1}(\tol) \le \rollmin_{\param_i}^{-1}(k) \\&\iff \rollmin_{\param_i}(\truegap_{\param_i}^{-1}(\tol)) \ge k,
    \end{align*}
    where the first equivalence is the definition of $\pst_k$ and the second is by~\eqref{eq:confproof1}.
    % \bnote{do we ever use $N_k$ except for this one time? If not, I would remove the definition}
    % Therefore, if we let 
    % $$
    % N_k = \cardinality\{i \in \{1, \dots, \ncalib\} \mid \rollmin_{\param_i}(\truegap_{\param_i}^{-1}(\tol)) \ge k \},
    % $$
    % by definition,
    Therefore, by definition, $\lt$ is exactly
    \begin{equation*}
        % \lt
        % &= \sup \{ k \ge 0 \mid  \ge n\} \\
        % &= \sup\{k \ge 0 \mid \cardinality\{i \in \{1, \dots, \ncalib\} \mid \truegap_{\param_i}^{-1}(\tol) \le \pst_k(\param_i) \} \ge n \} \\
        % &= \sup\{k \ge 0 \mid \cardinality\{i \in \{1, \dots, \ncalib\} \mid \truegap_{\param_i}^{-1}(\tol) \le \rollmin_{\param_i}^{-1}(k) \} \ge n \} \\
        \sup\{k \ge 0 \mid \cardinality\{i \in \{1, \dots, \ncalib\} \mid \rollmin_{\param_i}(\truegap_{\param_i}^{-1}(\tol)) \ge k \} \ge n \},
    \end{equation*}
    which corresponds to the $n$-th largest element of $\{\rollmin_{\param_i}(\truegap_{\param_i}^{-1}(\tol))\}_{i=1}^\ncalib$.
    By Lemma~\ref{lemma:ordering} applied to $Z_i = \rollmin_{\param_i}(\truegap^{-1}_{\param_i}(\tol))$,
    \begin{align*}
    \Prob[\rollmin_{\param_{\ncalib+1}}(\truegap^{-1}_{\param_{\ncalib + 1}}(\tol)) \ge \lt] &= \Prob[Z_{\ncalib + 1} \ge Z_{[\ncalib + 1 - n]}] \\ & \ge n / (\ncalib + 1).
    \end{align*}
    Reversing the equivalences above yields $\Prob[\truegap_{\param_{\ncalib+1}}(\predgap^{-1}_{\param_{\ncalib+1}}(\lt)) \le \tol] \ge n/(\ncalib + 1)$.
\end{proof}

To apply the theorem, fix tolerance $\tol > 0$ and error probability $\errorprob > 0$.
Given a \emph{calibration dataset} $\calibration = \{\param_i\}_{i=1}^\ncalib$ of independent draws from $\paramdist$, choose $n$ such that $n/(\ncalib+1) \ge 1 - \errorprob$.
Theorem~\ref{thm:conformal} yields a threshold $\lt$ such that for any new instance $\param_{\ncalib+1} \sim \paramdist$ independent of $\calibration$, terminating at time $\pst_{\lt}(\param_{\ncalib+1})$ gives
\begin{equation*}
    \Prob[\truegap_{\param_{\ncalib+1}}(\pst_{\lt}(\param_{\ncalib+1})) \le \tol] \ge 1 - \errorprob.
\end{equation*}

\change{
\paragraph{Remark (multiple testing).}
A naive approach would predict the optimality gap at every solver callback and terminate the first time the prediction drops below $\tol$.
This constitutes multiple testing: each callback represents a new hypothesis test, and the probability of at least one false positive (stopping before a $\tol$-optimal solution has been found) grows with the number of callbacks, which can number in the hundreds per solve.
Our conformal approach avoids this issue because the threshold $\lt$ is calibrated \emph{once} on the calibration set, prior to deployment.
The guarantee in Theorem~\ref{thm:conformal} covers the \emph{entire} trajectory: the single pre-calibrated threshold $\lt$ ensures that the first time $\predgap_\param(t) \le \lt$ occurs, the solution is $\tol$-optimal with the stated probability, with no multiple-testing correction required.
}

\subsection{Training the predictor} \label{sec:training}
We parametrize $\truegap_\param(t)$ using a neural network.
Let $\covars_\param(t) \in \reals^M$ be the solver state vector at time $t$, containing the upper and lower bounds on $\ov(\param)$, the current best feasible solution $\vars_\param(t)$, the number of explored nodes, the parameter $\param$, and the elapsed solve time.
Let $\nn: \reals^M \rightarrow \reals_+$ be a neural network.
For $l \le u$ define the squashing function
\begin{equation}
    \activ(x \mid l, u) = (u - l)\frac{e^x}{1 + e^x},
\end{equation}
which satisfies $\activ(x \mid l, u) \in [0, u-l]$ for all $x \in \reals$.
The predictor is
\begin{equation*}
\predgap_\param(t) = \activ\big(\nn(\covars_\param(t)) \mid \LB_\param(t), \UB_\param(t)\big).
\end{equation*}
By Assumption~\ref{ass:terminal}, $\UB_\param(t) - \LB_\param(t) \to 0$ as $t \to \infty$, so $\predgap_\param(t) \to 0$.

We train the predictor to approximate $\UB_\param(t) - \ov(\param)$ using the weighted loss
\begin{equation*}
    \loss = \E\left[ \int_0^\infty \weight_{\param}(t) \Big(\predgap_\param(t) - \big(\UB_\param(t) - \ov(\param)\big)\Big)^2 dt\right],
\end{equation*}
where $\weight_\param(t) \ge 0$ is a time-weighting with $\int_0^\infty \weight_{\param}(t) dt = 1$.
Let $T_\param$ denote the termination time of branch-and-bound on problem~\eqref{eq:master}.
For discretization step $\delta > 0$, let $S_\param = \lceil T_\param / \delta \rceil$.
The empirical loss over training data $\train = \{\param_i\}_{i=1}^\ntrain$ is
\begin{equation} \label{eq:loss2}
    \hat{\loss} = \frac{1}{\ntrain} \sum_{i=1}^\ntrain \sum_{t=0}^{S_{\param_i}} \weight_{\param_i}(t \delta) \Big(\predgap_{\param_i}(t \delta) - y_{\param_i}(t \delta)\Big)^2,
\end{equation}
where $y_{\param_i}(t \delta) = \UB_{\param_i}(t \delta) - \ov({\param_i})$.
We minimize~\eqref{eq:loss2} using stochastic gradient descent.
The weighting
\begin{equation}
    \weight_{\param}(t) \propto \frac{1}{y_{\param}(t)} \label{eq:weight}
\end{equation}
assigns highest weight to times when $y_{\param}(t)$ is small.
Prediction accuracy matters most near $y_{\param}(t) \approx \tol$, since this is when the termination decision is made.
Because $\tol$ is typically small, the weighting~\eqref{eq:weight} focuses the predictor on the critical region where $\truegap_{\param}(t) \approx \tol$.

% \subsection{Warm-starts and other heuristics}
% The conformal prediction method can be combined with warm-starting heuristics and any other heuristics which might be used to speed up the branch-and-bound algorithm, including learned heuristics~\cite{Shen2021LearningPH, Manchanda2022OnTG}.
% In our experiments section we fit the stopping time on models which have access to all optimal solutions from the training data as warm-starts (when appropriate) and also on models which have no extra information.

\subsection{Method overview} \label{sec:overview}
The overall pipeline comprises three steps.
\begin{enumerate}
    \item \textbf{Training.} Sample $\ntrain$ instances $\param \sim \paramdist$ to form $\train = \{\param_i\}_{i=1}^\ntrain$ and train the predictor $\truegap_\param(t)$ by minimizing~\eqref{eq:loss2}.
    \item \textbf{Calibration.} Sample $\ncalib$ instances $\param \sim \paramdist$ to form $\calibration = \{\param_i\}_{i=1}^\ncalib$.
    Fix tolerance $\tol > 0$ and error probability $\errorprob \in (0,1)$.
    Compute $\lt$ via Theorem~\ref{thm:conformal}.
    \item \textbf{Deployment.} On a new instance $\param$, terminate at time $\pst_\lt(\param)$ (when $\truegap_\param(t)$ first drops below $\lt$) and return the best feasible solution.
\end{enumerate}
The returned solution is $\tol$-optimal with probability at least $1 - \errorprob$ over the calibration set and test instance.
Figure~\ref{fig:method_overview} illustrates the method.

% \bnote{ I thought about the guarantees and I think the right way to mention them is sample-conditional \url{https://web.stanford.edu/class/stats315a/Reading/conformal-thoughts.pdf}}
\section{Theoretical guarantees}
Let the suboptimality at termination be
\begin{equation*}
    \termov(\param, \lt) = \frac{\change{\UB_\param(\pst_{\lt}(\param))} - \ov(\param)}{|\ov({\param})|}.
\end{equation*}
By construction, our method satisfies
\begin{equation}
\Prob[\termov(\param, \lt) \le \tol] \ge 1 - \errorprob, \label{eq:confprob2}
\end{equation}
where randomness is over the calibration set $\calibration$ and $\param \sim \paramdist$.
Define the empirical and expected suboptimality of threshold $\lt$ as
\begin{equation*}
    \hat{S}(\lt) = \frac{1}{\ncalib} \sum_{i=1}^\ncalib \termov(\param_i, \lt), \quad S(\lt) = \E_{\param \sim \paramdist} [\termov(\param, \lt) \mid \calibration],
\end{equation*}
and the empirical and expected runtime as
\begin{equation*}
    \hat{\runtime}(\lt) = \frac{1}{\ncalib} \sum_{i=1}^\ncalib \pst_\lt(\param_i), \quad \runtime(\lt) = \E_{\param \sim \paramdist} [\pst_\lt(\param) \mid \calibration].
\end{equation*}
The guarantee~\eqref{eq:confprob2} averages over both the calibration set and test point.
This section establishes stronger \emph{sample-conditional} guarantees~\cite{Duchi2025AFO} that hold for a fixed calibration set with high probability:
\begin{enumerate}
    \item \textbf{Expected suboptimality and runtime} (Section~\ref{sec:expected}, Theorem~\ref{thm:expect}): $S(\lt) \le \hat{S}(\lt) + O(1/\sqrt{\ncalib})$ and $\runtime(\lt) \le \hat{\runtime}(\lt) + O(1/\sqrt{\ncalib})$, so the expected suboptimality and runtime on new instances are close to their empirical values on calibration data.
    \item \textbf{Success probability} (Section~\ref{sec:success}, Theorem~\ref{thm:condprob}): $\Prob[\termov(\param, \lt) \le \tol \mid \calibration] \ge 1 - \alpha - O(1/\sqrt{\ncalib})$, so the conditional probability of achieving tolerance $\tol$ concentrates near $1 - \alpha$.
\end{enumerate}

% In this section, we will show that on top of the guarantee of Equation~\eqref{eq:confprob2}, our method satisfies the following guarantees on the quality of solutions returned on test points, often referred to as \emph{sample-conditional} guarantees~\cite{Duchi2025AFO}:
% \begin{enumerate}
%     \item For any $\delta > 0$ With probability at least $1 - \delta$ over the calibration set $\calibration \sim \paramdist^\ncalib$, conditional on the calibration data we have, $\E[z(\param,\lt)|\calibration] \le \epsilon + o(\sqrt{\ncalib}) \sqrt{\log(1/\delta)}$ and $\Prob[z(\param, \lt)\le\epsilon|\calibration] \ge 1 - \alpha - o(\sqrt{\ncalib}) \sqrt{\log(1/\delta)}$.
%     \item For any $\delta > 0$ With probability at least $1 - \delta$ over the calibration set $\calibration \sim \paramdist^\ncalib$, conditional on the calibration data we have, $\runtime(\lt) \le \hat{\runtime}(\lt)+ o(\sqrt{\ncalib}) \sqrt{\log(1/\delta)}$.
% \end{enumerate}

\subsection{Expected suboptimality and runtime} \label{sec:expected}
We bound the expected suboptimality and runtime conditional on the calibration set.
Let $\zbd$ bound the suboptimality, \ie, $\termov(\param, \lt) \le \zbd$ for all $\lt \ge 0$ and $\param \in \paramspace$.
Let $\tbd$ bound the runtime, \ie, $\pst_\lt(\param) \le \tbd$ for all $\lt \ge 0$ and $\param \in \paramspace$.
These bounds can be enforced via deterministic stopping criteria.
\begin{theorem} \label{thm:expect}
    For any $\delta > 0$, with probability at least $1 - \delta$ over the calibration set $\calibration$,
    \begin{equation*}
        S(\lt) \le \hat{S}(\lt) + \zbd \sqrt{\frac{2\log (e \ncalib)}{\ncalib}} + \zbd \sqrt{\frac{\log(1/\delta)}{2\ncalib}}.
    \end{equation*}
    The runtime satisfies the analogous bound
    \begin{equation*}
        \runtime(\lt) \le \hat{\runtime}(\lt) + \tbd \sqrt{\frac{2\log (e \ncalib)}{\ncalib}} + \tbd \sqrt{\frac{\log(1/\delta)}{2\ncalib}}.
    \end{equation*}
\end{theorem}
% We first prove a lemma which we need to control the VC-dimension of our regret.
% \begin{lemma} \label{lemma:2d}
%     Let $B_1, \dots B_n$ be sets in $\reals^2$ such that for each $(x,y) \in B_i$, if $\tilde{x} \ge x$ and $\tilde{y} \le y$ then also $(\tilde{x}, \tilde{y}) \in S_i$.
%     Suppose also that for every subset $S$ of $\{1, \dots n\}$ there is an $(x,y) \in \reals^2$ such that for all $i \in S$, $(x,y) \in B_i$ and for all $i \notin S$, $(x,y) \notin B_i$. 
%     Then $n$ is at most 2.
% \end{lemma}
\begin{proof}[Proof of Theorem~\ref{thm:expect}]
    We prove the suboptimality bound; the runtime bound follows analogously.
    The proof uses VC-dimension to obtain uniform convergence.
    For $u \in [0, \zbd]$ and $\lt \in \reals_+$ define the classifier
    \begin{equation*}
        \classif(\lt,u)(\param) = \indicator_{\termov(\param, \lt) > u},
    \end{equation*}
    which equals 1 if $\termov(\param, \lt) > u$ and 0 otherwise.
    Define the empirical and true risk as
    \begin{align*}
        \hat{S}(\classif(\lt, u)) =& \frac{1}{\ncalib}\sum_{i=1}^\ncalib \classif(\lt,u)(\param_i), \\ \quad S(\classif(\lt, u)) =& \underset{\param \sim \paramdist}{\E}[ \classif(\lt,u)(\param)].
    \end{align*}
    For fixed $u \in [0, \zbd]$ let $\vc(u)$ be the VC-dimension (Definition 3.10,~\cite{Mohri2012FoundationsOM}) of the set of functions $\{\classif(\lt, u) \mid \lt \in \reals_+\}$.
    Let $\hat{\paramdist}$ be a uniform distribution on~$\calibration$ (the empirical distribution for the calibration dataset).
    We have,
    \begin{align*}
        & |S(\lt) - \hat{S}(\lt)| \\ &= \Big|\underset{\param \sim \paramdist}{\E}[\termov(\param, \lt)] - \underset{\param \sim \hat{\paramdist}}{\E}[\termov(\param, \lt)]\Big| \\
        &= \Bigg| \int_{0}^\zbd \underset{\param \sim \paramdist}{\Prob}[\termov(\param, \lt)>u] - \underset{\param \sim \hat{\paramdist}}{\Prob}[\termov(\param, \lt)>u] \text{d}u \Bigg| \\
        &\le\zbd  \underset{u \in [0, \zbd]}{\sup} \big| S(\classif(\lt, u)) - \hat{S}(\classif(\lt, u)) \big| \\
        & \le \zbd \underset{u \in [0, \zbd]}{\sup}\sqrt{\frac{2 \vc(u) \log(e \ncalib /\vc(u))}{\ncalib}} + \zbd \sqrt{\frac{\log(1/\delta)}{2\ncalib}}.
    \end{align*}
    The first equality is the definition of the empirical distribution.
    The second uses $\E[X] = \int_0^\infty \Prob[X > u] \, du$ for nonnegative $X$.
    The third bounds the integral uniformly.
    The fourth applies Corollary 3.19 of~\citet{Mohri2012FoundationsOM}.
    It remains to show $\vc(u) = 1$.
    The function $\lt \mapsto \classif(\lt, u)(\param)$ is nondecreasing in $\lt$.
    For any $\param_1, \param_2$, let $\lt_i = \inf\{\lt : \classif(\lt, u)(\param_i) = 1\}$.
    If $\lt_1 \le \lt_2$, then $\classif(\lt, u)(\param_1) \ge \classif(\lt, u)(\param_2)$ for all $\lt$, so $\{\param_1, \param_2\}$ cannot be shattered.
    Thus $\vc(u) \le 1$, and since $\vc(u) \ge 1$ by definition, we have $\vc(u) = 1$.
    % Since $\termov$ is nondecreasing as a function of $\lt$, for a given $\param$, the function $\classif(\lt, u)(\param)$ must be increasing in $\lt$ and decreasing in $u$.
    % For a given $\param$, define the set,
    % $A(\param) = \{(\lt, u) | \rho(\lt, u)(\param) = 1\}$.
    % The definition of VC dimension is the maximum integer $n$ we can find such that there exists $\param_1, \dots \param_n \in \paramspace$ such that for every subset $S$ of $\{1, \dots n\}$, there is an $(\lt, u) \in \reals^2$ such that $i \in S \iff (\lt, u) \in A(\param_i)$. 
    % Each set $A(\param)$ also satisfies the two-dimensional monotonicity property that $(\lt, u) \in A(\param)$, $\tilde{\lt} \ge \lt, \tilde{u} \le u \implies (\tilde{\lt}, \tilde{u}) \in A(\param)$. 
    % Therefore the statement is true by Lemma~\ref{lemma:2d}.
\end{proof}

\subsection{Success probability} \label{sec:success}
The probability in~\eqref{eq:confprob2} is over both the calibration set and test point, which is equivalent to
\begin{equation*}
    \E_{\calibration \sim \paramdist^\ncalib}\Big[ \Prob_{\param \sim \paramdist}[\termov(\param, \lt) \le \tol \mid \calibration]\Big] \ge 1 - \alpha.
\end{equation*}
The next result shows that the conditional probability $\Prob[\termov(\param, \lt) \le \tol \mid \calibration]$ concentrates near $1 - \alpha$.
The proof follows from Proposition 2a of~\citet{Vovk2012ConditionalVO}; we include it for completeness.
\begin{theorem} \label{thm:condprob}
    For any $\delta > 0$, with probability at least $1 - \delta$ over the calibration set $\calibration$,
    \begin{equation*}
        \underset{\param \sim \paramdist}{\Prob}[\termov(\param, \lt) \le \tol \mid \calibration] \ge 1 - \alpha - \sqrt{\frac{\log(2/\delta)}{2\ncalib}}.
    \end{equation*}
\end{theorem}
\begin{proof}
    Let $\hat{\paramdist}$ be the uniform distribution on $\calibration$ (the distribution over the empirical sample).
    By definition of $\lt$ we know that,
    \begin{equation} \label{eq:proof1}
        \underset{\param \sim \hat{\paramdist}}{\Prob} [\termov(\param, \lt) \le \tol] \ge 1 - \alpha.
    \end{equation}
    Let $\hat{P}_\epsilon = \underset{\param \sim \hat{\paramdist}}{\Prob} [\termov(\param, \lt) \le \tol]$ and $P_\epsilon = \underset{\param \sim\paramdist}{\Prob} [\termov(\param, \lt) \le \tol]$.
    By Corollary 1~\cite{Massart1990TheTC} the following is true, for all $\lambda > 0$, where the randomness is taken over the calibration set $\calibration$,
    \begin{equation*}
        \Prob[|\hat{P}_\epsilon -  P_\epsilon| \ge \lambda /\sqrt{\ncalib}] \le 2 \exp(-2 \lambda^2).
    \end{equation*}
    Set $\delta = 2 \exp(-2 \lambda^2)$ and rearrange this expression to get,
    \begin{equation*}
    \Prob\Bigg[|\hat{P}_\epsilon -  P_\epsilon| \ge \sqrt{\frac{\log(2/\delta)}{2\ncalib}}\Bigg] \le \delta.
    \end{equation*}
    This combined with~\eqref{eq:proof1} completes the proof of the first inequality.
\end{proof}
\subsection{Robustness to distribution shifts}
\change{
    We can also show that the conformal guarantees hold under a shifted distribution $\tilde{\paramdist}$ which is $\rho$-close to $\paramdist$ in total variation distance, so the method is robust to distribution shifts.
}
\begin{theorem} \label{thm:distshift}
    \change{
        Let $\tilde{\paramdist}$ be a distribution such that $\operatorname{TV}(\tilde{\paramdist}, \paramdist) \le \rho$.
        Then the guarantees of Theorems~\ref{thm:expect} and~\ref{thm:condprob} hold with respect to $\tilde{\paramdist}$, with an additional additive $\rho$ term.
        In particular, we have,
        \begin{equation*}
            \Prob[\termov(\param, \lt) \le \tol] \ge 1 - \alpha - \rho,
        \end{equation*}
        where the probability is over the shifted distribution test point $\param \sim \tilde{\paramdist}$ and calibration set $\calibration \sim \paramdist^\ncalib$.
        Further, we have, with probability at least $1 - \delta$ over the calibration set $\calibration$,
        \begin{equation*}
            \underset{\param \sim \tilde{\paramdist}}{\Prob}[\termov(\param, \lt) \le \tol \mid \calibration] \ge 1 - \alpha - \rho - \sqrt{\frac{\log(2/\delta)}{2\ncalib}}.
        \end{equation*}
    }
\end{theorem}
\proof
    \change{
        Write $\lt(\calibration)$ for the random variable $\lt$ as a function of the calibration data $\calibration$.
        For the first inequality,
        \begin{align*}
            & \Prob\left[\termov(\param, \lt) \le \tol\right] \\ &= \underset{\calibration}{\E}\left[\underset{\param \sim \tilde{\paramdist}}{\Prob}\left[\termov(\param, \lt(\calibration)) \le \tol\right]\right] \\
            &= \underset{\calibration}{\E}\left[\underset{\param \sim \paramdist}{\Prob}\left[\termov(\param, \lt(\calibration)) \le \tol\right] \right. \\ & \qquad \left. - \left| \underset{\param \sim \tilde{\paramdist}}{\Prob}\left[\termov(\param, \lt(\calibration)) \le \tol\right] -  \underset{\param \sim \paramdist}{\Prob}\left[\termov(\param, \lt(\calibration)) \le \tol\right] \right| \right] \\
            & \ge \underset{\calibration}{\E}\left[\underset{\param \sim \paramdist}{\Prob}\left[\termov(\param, \lt(\calibration)) \le \tol\right]\right] - \rho \\
            &\ge 1 - \alpha - \rho.
        \end{align*}
        For the second inequality let $$\tilde{P} = \underset{\param \sim \tilde{\paramdist}}{\Prob}[\param \in \{\param \mid \termov(\param, \lt) \le \tol\} \mid \calibration]$$ and $P = \underset{\param \sim \paramdist}{\Prob}[\param \in \{\param \mid \termov(\param, \lt) \le \tol\} \mid \calibration]$. Then we have,
        \begin{align*}
            & \underset{\param \sim \tilde{\paramdist}}{\Prob}[\termov(\param, \lt) \le \tol \mid \calibration] \\ &= \underset{\param \sim \tilde{\paramdist}}{\Prob}[\param \in \{\param \mid \termov(\param, \lt) \le \tol\} \mid \calibration]
            \\ &= \underset{\param \sim \paramdist}{\Prob}[\param \in \{\param \mid \termov(\param, \lt) \le \tol\} \mid \calibration] - \Big| \tilde{P} - P \Big|
            \\ &\ge \underset{\param \sim \paramdist}{\Prob}[\termov(\param, \lt) \le \tol \mid \calibration] - \rho \\
            &\ge 1 - \alpha - \rho - \sqrt{\frac{\log(2/\delta)}{2\ncalib}}.
        \end{align*}
    }
\qed
\subsection{Time to termination}
The guarantees above hold for any predictor, but meaningful speedup requires an accurate predictor $\predgap_\param(t)$.
Taking $\lt = -\infty$ achieves zero suboptimality but provides no speedup.
The threshold $\lt$ from Theorem~\ref{thm:conformal} satisfies all guarantees regardless of predictor quality, but if $\predgap_\param(t)$ is inaccurate, the stopping time $\pst_\lt(\param)$ may equal the standard termination time $\st_\tol(\param)$.
\section{Computational experiments}
We test our methods on parametric families of problems from the Distributional MIPLIB library~\cite{Huang2024DistributionalMA}.
\change{We include a series of plots and figures demonstrating the statistical variety of the datasets considered in Appendix~\ref{appendix:experimentplots}.}
\subsection{Experiment setup}
\paragraph{Hardware.} All models were trained on a NVIDIA A100 GPU.
Each optimization solve took place on AMD EPYC 9334 CPUs, and each optimization solve was allowed access to exactly 1GB of memory, and only one thread.
\paragraph{Predictor.} We parametrize $\nn$ by an LSTM with $\LSTMLAYERS$ layers and $\LSTMHIDDEN$ hidden units in each layer, followed by a feedforward neural network with $\FFLAYERS$ hidden layers and $\FFHIDDEN$ units in each layer.
All activation units but the final one are ReLU units.
The output dimension is one, and the output represents the approximate true optimality gap at time $t$.
We have $\predgap_\param(t) = \activ\big(\nn(\covars_\param(t)) \mid \LB_\param(t), \UB_\param(t)\big)$.
\change{
    We choose an LSTM because it is a simple model which can detect complex patterns in time-series data, but this is not a necessary part of our work. This learned model can be swapped for the most appropriate model for a given task.
    In some applications, the optimization problem might be sufficiently hard, or the solve might take sufficiently long, that a simple RNN/LSTM might not perform well. 
    In this case, practitioners may want to swap this component out for a transformer, or include a graph neural network to add awareness of the entire problem formulation.
}
\paragraph{Covariates.} At time $t$, the covariate vector $\covars_\param(t)$ consists of the current bounds $\UB_\param(t)$ and $\LB_\param(t)$, their 1, 3, and 5 second rolling averages, the elapsed solve time $t$, and the number of nodes expanded in the branch-and-bound tree.
Unlike typical learn-to-optimize settings~\cite{Scavuzzo2022LearningTB}, the predictor does not require access to the full problem formulation, so generating covariate data is efficient.
%The part of the process which uses the most computational resources is the collection of solver data and covariates for each problem instance in the training and calibration datasets in step 1 of Section~\ref{sec:overview}.

\begin{table}
    \centering
    \caption{The number of training, calibration, and test instances for each problem family, along with the corresponding reference. \change{The OTS-medium family uses $\ncalib = 28$ and $\ntest = 72$ because only 100 instances are available in total from the real-world dataset.}}
    \label{tab:problems}
    \small
\begin{tabular}{@{}l l l l l@{}}
\toprule
Problem & $\ntrain$ & $\ncalib$ & $\ntest$ & Reference \\
\midrule
CFLP-medium & 800 & 100 & 100 & \citeauthor{Scavuzzo2022LearningTB} \\
GISP-easy & 800 & 100 & 100 & \citeauthor{hochbaum1997forest} \\
OTS-medium & 800 & 28 & 72 & \citeauthor{Huang2024DistributionalMA} \\
MMCN-medium-BI & 800 & 100 & 100 & \citeauthor{greening2023lead} \\
MIS-medium & 800 & 100 & 100 & \citeauthor{Huang2024DistributionalMA} \\
MVC-medium & 800 & 100 & 100 & \citeauthor{korte2018combinatorial} \\
\bottomrule
\end{tabular}
\end{table}

\paragraph{Optimization solver.}
The solvers we test our method on are \change{Gurobi 12}~\cite{gurobi} and Cardinal Optimization~\cite{Ge2022CardinalO}.
We use solver callbacks to evaluate the predictor $\predgap_\param(t)$ as the solve is running.
We disable presolve on both solvers, use only a single thread, \change{and set Gurobi MIPFocus parameter to 1}. 
The maximum allowed solve time is set to 600 seconds.
In our experiments we seek to solve all problem instances to within $0.1\%$ accuracy, so we set $\epsilon = 0.001$.
We choose $\alpha = 0.05$, so that our solver is likely to return an $\tol$-optimal solution at least $95\%$ of the time.
% \paragraph{Data.}
% We consider 5 families from the standard distributional MIPLIB library of parametric families of MIPs.
% Each problem family is split into a training dataset of $\ntrain$ problems, a calibration dataset of $\ncalib$ problems, and a test dataset of $\ntest$ problems.
% The number of problems in each dataset varies depending on the problem family.
% For training, we produce data on $\covars_\param(t)$ for each instance $\param$ by solving the instance with Gurobi and saving the solver data with callbacks.
% This is used to train the predictor.
% The calibration set is used to calibrate the conformal predictor, as described in~\eqref{sec:conformal}.
% We display results on the test datasets.
\paragraph{Evaluation metrics.}
For each problem family, we report average suboptimality on test data $\test = \{\param_i\}_{i=1}^\ntest$,
\begin{equation*}
    \hat{S} = \frac{1}{\ntest}\sum_{i=1}^\ntest \termov(\param_i, \lt),
\end{equation*}
where $\termov(\param_i, \lt)$ is the relative suboptimality of the solution returned by our method on instance $\param_i$, as defined in Section~\ref{sec:expected}.
% We refer to the first as the absolute test regret and the second as the relative test regret.
We also report total solve time and the percent of problems solved to within $\tol$-optimality.
\paragraph{Baselines.}
We compare our method to the baseline methods in Table~\ref{tab:baselines}.
\change{Comparisons to a wider range of baselines are also included in Appendix~\ref{appendix:moreexperiments}.}
\begin{table}[h]
	\centering
	\caption{Baseline methods.}
	\label{tab:baselines}
	\begin{tabular}{lll}
		\toprule
		Method & Solver & Termination criterion \\
		\midrule
		GRB & Gurobi & $\tol$ optimality \\
		GRB1 & Gurobi & 1 solution found \\
		GRB3 & Gurobi & 3 solutions found \\
		COPT & COPT & $\tol$ optimality \\
		COPT1 & COPT & 1 solution found \\
		COPT3 & COPT & 3 solutions found \\
		\bottomrule
	\end{tabular}
\end{table}

\paragraph{Data.}
We run experiments on capacitated facility location problems (CFLP-medium), generalized independent set problems (GISP-easy), maximum independent set problems (MIS-medium), minimum vertex cover problems (MVC-medium), optimal transmission switching problems (OTS-medium), and middle-mile consolidation network design problems (MMCN-medium-BI).
The OTS-medium and MMCN-medium-BI problems come from real-world datasets.
 These problems are costly to solve, and on some families, Gurobi takes hundreds of seconds to solve most instances.
The number of train points $\ntrain$, calibration points $\ncalib$, and test points $\ntest$ vary depending on the problem family.
The details on data and formulations are given in Table~\ref{tab:problems}.

% \begin{itemize}
%     \item Capacitated facility location problems (CFLP-medium): This is a synthetic dataset for a problem which is commonly used to benchmark MIP solvers. The full formulation is given in~\citet{Scavuzzo2022LearningTB}.
%     Here $\ntrain = 800$, $\ncalib = 100$, and $\ntest = 100$.
%     \item Generalized independent set problems (GISP-easy):  Another common synthetic benchmark for MIPs.
%     The formulation is given in~\citet{hochbaum1997gisp}. In this dataset $\ntrain=800$, $\ncalib = 100$, and $\ntest = 100$.
%     \item Maximum independent set problems (MIS-medium): This is a synthetic dataset for a problem which is commonly used to benchmark MIP solvers. In this dataset $\ntrain=800$, $\ncalib = 100$, and $\ntest = 100$.
%     The full formulation is given on the distributional MIPLIB website~\cite{Huang2024DistributionalMA}.
%     \item Minimum vertex cover problems (MVC-medium): A problem where the goal is to find the smallest set of vertices in a graph such that each edge is incident to at least one vertex in the set. This is another synthetic problem.
%     The full formulation is available in~\cite{korte2018combinatorial}. In this dataset $\ntrain=800$, $\ncalib = 100$, and $\ntest = 100$.
%     \item Optimal transmission switching problems (OTS-medium): This is a real-world dataset with data from a real power grid.
%     The original formulation is given in~\cite{pollack2024equitably}.
%     Here $\ntrain = 800$, $\ncalib = 72$ and $\ntest = 28$.
%     The aim of the problem is to either de-energize or to move underground a set of power lines in an electrical grid to minimize the risk of wildfire while minimizing power outages within a budget.
%     \item Neural network verification problems (NNV-easy): The data for these problems were generated originally by Google Deepmind~\cite{Nair2020SolvingMI}.
%     It contains a set of mixed integer programs (MIPs) for the problem of verifying a neural network's robustness to perturbations to its inputs.
%     Here $\ntrain = 2555$, $\ncalib = 550$, and $\ntest = 588$.
%     The data for this problem comes from neural networks trained to classify MNIST digits, so is not synthetic.
% \item Middle-mile consolidation network design problems (MMCNDP): 
% In this problem a distribution company needs to move shipments from a set of origins to a set of destinations, within time constraints.
% The company has a fleet of vehicles with limited capacity and aims to design a network that minimizes the total transportation cost while satisfying all demand within the given time windows.
% The original formulation is given in~\citet{greening2023lead}.
% This is a real-world dataset using supply-chain data.
% Here $\ntrain = 800$, $\ncalib = 100$, and $\ntest = 100$.

\begin{figure}[t]
    \centering
    \ifarxiv
    \includegraphics[width=0.8\columnwidth]{solved_boxplots.pdf}%
    \else
    \includegraphics[width=\columnwidth]{solved_boxplots.pdf}%
    \fi
    \caption{Suboptimality, solve time, and number of nodes explored at termination across all problem families.
        The dashed line in the top panel marks the target suboptimality $\tol = 0.1\%$; values are clipped to $10^{-6}$ when an optimal solution is found.
        Our methods (shaded grey) achieve substantially faster solve times than the global solvers Gurobi (GRB) and COPT while maintaining low suboptimality.}
    \label{fig:results}
\end{figure}

\begin{figure}[t]
    \centering
    \ifarxiv
    \includegraphics[width=0.8\columnwidth]{solved_plot.pdf}%
    \else
    \includegraphics[width=\columnwidth]{solved_plot.pdf}%
    \fi
    \caption{
        The number of instances solved to within $2\epsilon = 0.2\%$-optimality over time for each method on each dataset.
        % Here Gurobi refers to Gurobi solving to $\tol$-optimality.
        Since $\epsilon$ was the conformal target tolerance, we expect almost all problems from each method to be solved to this level of optimality.
        The conformal prediction method solves most problems very quickly, but it is not guaranteed to solve all problems to within $0.2\%$-optimality due to the adjustment made by Theorem~\ref{thm:conformal}, which occurs with probability around $\errorprob$.
        % The grey dotted line marks $(1 - \alpha)$.
    }
    \label{fig:solvetime2}
\end{figure}


% \end{itemize}
\paragraph{Code overview.}
We first solve all problems in the training dataset $\{\param_i\}_{i=1}^{\ntrain}$ using solver callbacks and save the solver state $\covars_t(\param_i)$ at each point $t$ and true optimality gap $\truegap_{\param_i}(t)$ at which the callback is called.
This process can be computationally intensive. 
The next step is to train the predictor $\predgap_\param(t)$ using the saved solver states and true optimality gaps.
We run the Adam optimizer~\cite{Kingma2015Adam} using the PyTorch~\cite{Paszke2019PyTorchAI} automatic differentiation software to minimize the loss~\eqref{eq:loss2} where $\predgap_\param(t) = \activ\big(\nn(\covars_\param(t)) \mid \LB_\param(t), \UB_\param(t)\big)$.
The training process runs for half an hour on a single GPU for each experiment.
We then perform the conformal prediction step described in Section~\ref{sec:conformal} to get $\lt$, and evaluate the quality of the stopping time $\pst_\lt(\param)$ on the test dataset.
Code for our experiments is available through the following link:
\begin{center}
    \href{https://github.com/stellatogrp/conformal_mip}{\nolinkurl{github.com/stellatogrp/conformal_mip}}
\end{center}

% \subsection{Capacitated facility location problems}
% We first consider the `CFLP-medium' dataset from Distributional MIPLIB~\cite{Huang2024DistributionalMA}.
% This is a synthetic dataset for a problem which is commonly used to benchmark MIP solvers.
% \paragraph{Problem.}
% In this problem, $I$ is a set of clients and $J$ is a set of sites where facilities can be located.
% The objective is to match all clients to facilities in a way that minimizes the total cost.
% The variable $\faclocx_j \in \{0,1\}$ indicates whether the facility at location $j \in J$ is built, and the variable $\faclocy_{i,j} \in [0,1]$ represents the proportion of client $i$'s demand serviced by facility $j$.
% The cost of building a facility at location $j$ is $c_j(\param) \in \reals_+$, and the cost of servicing customer $i$ from location $j$ is $d_{ij}(\param) \in \reals_+$.
% The demand for customer $i$ is given by $e_{i}(\param) \in \reals_+$, and the capacity of location $j$ is given by $s_j \in \reals_+$.
% The problem is formulated as
% \begin{equation*} \label{eq:facloc}
%     \begin{array}{ll}
%         \text{minimize} & \sum_{i \in I, j \in J} d_{ij}(\param) \faclocy_{ij} + \sum_{j \in J} c_j(\param) \vars_j \\
%         \text{subject to} 
%         & \sum_{i \in I} e_i(\param) \faclocy_{i j} \le s_j \vars_j, \quad  j \in J, \\
%         & \faclocy_{ij} \in [0, 1], \quad i \in I, j \in J, \\ & \vars_j \in \{0, 1\}, \quad i \in I, j \in J.
%     \end{array}
% \end{equation*}


% \paragraph{Results.} Results are displayed in Figure~\ref{fig:results}.

% \subsection{Generalized independent set problems}
% We consider the `GISP-easy' dataset from Distributional MIPLIB~\cite{Huang2024DistributionalMA}.
% This is another synthetic dataset, but the GISP problem is standard and commonly used in benchmarking.
% \paragraph{Problem.}
% The original problem formulation is defined in~\citet{hochbaum1997forest} and~\citet{Colombi2017TheGI}.
% In this problem, we are given a graph $G = (V, E)$ with vertex set $V$ and edge set $E$.
% The set of edges are partitioned into a set of non-removable edges $E_1$ and a set of removable edges $E_2$.
% A revenue $w_v(\param)$ is associated with each vertex $v \in V$ and a cost $c_e(\param)$ is associated with each removable edge $e \in E_2$.
% The goal of the optimization problem is to find a subset of vertices $S \subseteq V$ and a subset of removable edges $F \subseteq E_2$ such that no two edges in the vertex set are adjacent which
% maximizes the total revenue of the vertices minus the cost of the removed edges.
% Edges may only be removed from the set of removable edges $E_2$.
% Overall, the problem can be formulated as,
% \begin{equation*}
%     \begin{array}{ll}
%         \text{maximize} & \sum_{v \in V} w_v(\param) x_v - \sum_{e \in E_2} c_e(\param) z_e \\
%         \text{subject to} & \vars_u + \vars_v \le 1, \quad (u,v) \in E_1, \\
%         & x_u + x_v \le 1 + z_{(u,v)}, \quad (u,v) \in E_2, \\
%         & x_v \in \{0, 1\}, \quad v \in V, \\
%         & z_e \in \{0, 1\}, \quad e\in E_2.
%     \end{array}
% \end{equation*}
% \paragraph{Results.} Results are displayed in Figure~\ref{fig:results}.

% \subsection{Optimal transmission switching problems}
% We consider the `OTS-easy' dataset from Distributional MIPLIB~\cite{Huang2024DistributionalMA}.
% This is a subset of network topology optimization problems.
% The original formulation is given in~\cite{pollack2024equitably}.
% The aim of the problem is to either de-energize or to move underground a set of power lines in an electrical grid to minimize the risk of wildfirew chile minimizing power outages within a budget.
% This data comes from a real-world electrical grid.

% \paragraph{Problem.} In this problem, we 

% \subsection{Neural network verification problems}
% We consider the `NNV-easy' dataset from Distributional MIPLIB~\cite{Huang2024DistributionalMA}.
% This data was generated originally by Google Deepmind~\cite{Nair2020SolvingMI}.
% It contains a set of mixed integer programs (MIPs) for the problem of verifying a neural network's robustness to perturbations to its inputs.
% The data for this problem comes from neural networks trained to classify MNIST digits~\cite{Nair2020SolvingMI}, so is not synthetic.

% \paragraph{Problem.} We do not include a full problem formulation here, because the problems aer very complex and vary based on network architecture and parameters, but a rough formulation is the following,
% \begin{equation}
%     \begin{array}{ll}
%         \text{minimize} & \\
%         \text{subject to} &
%     \end{array}
% \end{equation}
% \paragraph{Results.} Results are displayed in Figure~\ref{fig:results}.

\subsection{Experiment results}
Figure~\ref{fig:results} shows suboptimality, solve time, and branch-and-bound nodes at termination for each method.
Full tabular results appear in Appendix~\ref{sec:tables}.
Across all problem families, the conformal prediction method maintains low suboptimality while substantially reducing solve time.
Compared to Gurobi, our method terminates 76\% faster on CFLP, 78\% faster on OTS, 66\% faster on MMCN, 88\% faster on MVC, and 95\% faster on MIS.
The improvement on GISP is more modest at 8\%, reflecting the structure of these instances: both solvers close the optimality gap relatively quickly after finding a good solution, leaving less room for early termination.

Figure~\ref{fig:solvetime2} shows the number of instances solved to $\tol$-optimality over time.
The conformal prediction method solves at least $100(1-\errorprob)\%$ of problems to $\tol$-optimality with high probability, though it may not solve all problems due to the probabilistic guarantee of Theorem~\ref{thm:conformal}.
Within the 600-second time limit, our method solves more instances to $\tol$-optimality than Gurobi alone on several problem families.
The improvement of CP~COPT over COPT is more modest but still clear.
This difference arises because Gurobi's heuristics find feasible solutions quickly, then spend most of the solve time tightening lower bounds to prove optimality.
Our predictor exploits this behavior by estimating when the current solution is already near-optimal, enabling early termination before the lower bound fully closes.


% \begin{figure}
%     \centering
%     \includegraphics[width=0.4\textwidth]{EXP_PLOT.pdf}%
%     \caption{Relative suboptimality and solve-time for each method on each dataset.
%     \snote{Deal with the rounding}
%     \snote{Do box chart instead of bar plot}
%     The conformal prediction method achieves low regret while significantly reducing solve-time across all datasets.
%     }
%     \label{fig:solvetime}
% \end{figure}
% \begin{figure}
%     \centering
%     \includegraphics[width=0.45\textwidth]{fig.pdf}%
%     \caption{Relative suboptimality versus solve-time for each method on each dataset.
%     Each point represents a single test instance.
%     The conformal prediction method achieves low regret while significantly reducing solve-time across all datasets.
%     }
%     \label{fig:time_vs_rel_subopt}
% \end{figure}


% \subsection{Middle-mile consolidation network design problems}
% The final dataset we consider is the `MMCND-Medium-BI' dataset from Distributional MIPLIB~\cite{Huang2024DistributionalMA}.
% The original formulation is given in~\citet{greening2023lead}.
% This is a real-world dataset using supply-chain data.

% \paragraph{Problem.} In this problem, we are given a set of potential warehouse locations $W$ and a set of retail locations $R$.

% \paragraph{Results.} Results are displayed in Figure~\ref{fig:results}.

\section{Conclusion}
Branch-and-bound solvers often find optimal solutions long before they can prove optimality.
We exploit this by learning a termination criterion that stops the solver early while providing rigorous probabilistic guarantees on solution quality.
Our method trains a neural network to predict the true optimality gap from solver state, then uses conformal prediction to calibrate a stopping threshold that ensures $\tol$-optimal solutions with probability at least $1 - \errorprob$.
We establish sample-conditional guarantees showing that expected suboptimality and success probability concentrate near their empirical values on calibration data.
On benchmarks from the Distributional MIPLIB library, our method achieves speedups exceeding 60\% on five of \change{six} problem families while maintaining $0.1\%$-optimality with 95\% probability.

\paragraph{Limitations.}
Our method requires solving instances from the target distribution during training to record solver state, which can be computationally expensive.
This means the approach accelerates problems that current solvers can already solve; it does not extend the frontier of tractable instances.
The speedup also varies across problem families: while most benchmarks show improvements exceeding 60\%, the GISP family achieves only 8\% speedup, suggesting that the gap between finding and proving optimality is problem-dependent.
Finally, the conformal guarantee assumes test instances are drawn from the same distribution as calibration data; distribution shift may degrade performance.


%     \caption{Tabular results for the experiments.
%     Time (s) refers to the average solve time over the test instances, plus or minus the standard deviation.
%     RelSubopt (\%) refers to the average relative suboptimality of the solutions over the test instances, plus or minus the standard deviation.
%     Correct (\%) refers to the percentage of instances for which the method found an $\tol$-optimal solution.
%     While there is little improvement on the smaller problem instances,
%     the conformal prediction method achieves low suboptimality while significantly reducing solve-time across all large problem instances.}
%     \label{fig:results}
% }


% \end{figure}


\newcommand{\myack}{%
Bartolomeo Stellato is supported by the NSF CAREER Award ECCS-2239771 and the ONR YIP Award N000142512147.
The authors are pleased to acknowledge that the work reported on in this paper was substantially performed using Princeton University's Research Computing resources.%
}

\ifarxiv
\section*{Acknowledgements}
\myack

\bibliographystyle{icml2026}
\bibliography{bibliography}
\else
\balance
\ifdefined\isaccepted
\section*{Acknowledgements}
\myack
\fi

\section*{Impact Statement}
This paper presents work on accelerating combinatorial optimization using probabilistic methods.
This has the potential to significantly reduce computational costs in real-world applications where combinatorial optimization problems are prevalent.

\bibliography{bibliography}
\bibliographystyle{icml2026}
\fi


\appendix
\onecolumn
\change{
\section{Experimental data} \label{appendix:experimentplots}
In this section we display plots and tables to demonstrate that the datasets and convergence curves for the problems considered are nontrivial.
\subsection{Convergence curve analysis}
Here we display information on the convergence curves of the MIPs.
Convergence curves for 20 instances of each test set are displayed in Figure~\ref{fig:convergencecurves}.
The following table numerically summarizes the convergence curve data.
Here UB conv. (s) is the average time taken (mean $\pm$ SD) for a run of Gurobi to find the optimal solution on the training dataset.
LB conv. (s) is the average time taken for a run of Gurobi to reduce the relaxation gap to zero in branch-and-bound.
Let $\tau^\star$ be the time at which the optimal solution is discovered by Gurobi. 
Certified gap at $\tau^\star$ is the relaxation gap defined by the branch-and-bound tree at $\tau^\star$, while Predicted gap at $\tau^\star$ is the gap estimated by our model $g_\param(\tau^\star)$.
Observe that our model reliably predicts that the gap is small at the time when the optimal solution is found.
Note that lines in the plot often start some time after zero. This is because lower-bounds only exist once the first node in the branch-and-bound tree has been evaluated and upper-bounds only exist once a feasible solution has been found.
\begin{table}[ht]
\centering
\caption{Convergence statistics across problem families (mean $\pm$ std).}
\label{tab:convergence_stats}
\begin{tabular}{lcccc}
\toprule
Family & UB conv.\ (s) & LB conv.\ (s) & Certified gap at $\tau^\star$ (\%) & Predicted gap at $\tau^\star$ (\%) \\
\midrule
CFLP  & $6.8 \pm 7.3$    & $15.6 \pm 40.8$   & $0.108 \pm 0.069$ & $0.034 \pm 0.022$ \\
GISP  & $19.0 \pm 13.0$  & $37.0 \pm 8.9$    & $20.3 \pm 10.3$   & $0.845 \pm 1.791$ \\
MIS   & $10.8 \pm 8.8$   & $221.0 \pm 209.4$ & $1.85 \pm 0.68$   & $0.014 \pm 0.037$ \\
MMCN  & $61.4 \pm 110.7$ & $245.2 \pm 230.7$ & $2.67 \pm 1.88$   & $0.149 \pm 0.177$ \\
MVC   & $9.7 \pm 7.8$    & $126.7 \pm 188.1$ & $0.70 \pm 0.18$   & $0.009 \pm 0.020$ \\
OTS   & $84.2 \pm 18.4$  & $431.8 \pm 139.7$ & $4.09 \pm 1.66$   & $0.018 \pm 0.071$ \\
\bottomrule
\end{tabular}
\end{table}

\begin{figure}[ht]
    \centering
    \begin{subfigure}{0.48\textwidth}
        \includegraphics[width=\linewidth]{convergence_curves_CFLP-medium_1779730813.pdf}
        \caption{CFLP-medium}
        \label{fig:cflpinfo}
    \end{subfigure}\hfill
    \begin{subfigure}{0.48\textwidth}
        \includegraphics[width=\linewidth]{convergence_curves_GISP-easy_1779730867.pdf}
        \caption{GISP-easy}
        \label{fig:gispinfo}
    \end{subfigure}

    \vspace{1em}
    \begin{subfigure}{0.48\textwidth}
        \includegraphics[width=\linewidth]{convergence_curves_MVC-medium_1779730943.pdf}
        \caption{MVC-medium}
        \label{fig:mvcinfo}
    \end{subfigure}\hfill
    \begin{subfigure}{0.48\textwidth}
        \includegraphics[width=\linewidth]{convergence_curves_OTS-medium_1779730892.pdf}
        \caption{OTS-medium}
        \label{fig:otsinfo}
    \end{subfigure}

    \vspace{1em}
    \begin{subfigure}{0.48\textwidth}
        \includegraphics[width=\linewidth]{convergence_curves_MMCN-medium-BI_1779730918.pdf}
        \caption{MMCN-medium-BI}
        \label{fig:mmcninfo}
    \end{subfigure}\hfill
    \begin{subfigure}{0.48\textwidth}
        \includegraphics[width=\linewidth]{convergence_curves_MIS-medium_1779730842.pdf}
        \caption{MIS-medium}
        \label{fig:misinfo}
    \end{subfigure}
    \caption{Convergence curves for 20 test instances from each problem family.}
    \label{fig:convergencecurves}
\end{figure}
%
\subsection{SCIP Heuristics}
We attempt to solve each MIP in each problem test-set with three SCIP~\cite{Hojny2025TheSO} heuristics: Feasibility Pump, RENS, and Simple Rounding.
We choose the SCIP solver for this task because, unlike the other solvers considered in this paper, it allows the user to enable and disable individual heuristics easily.
Results are displayed in Table~\ref{tab:scip_heuristics}.
Individual heuristics do not suffice to reliably solve problems in our dataset.
%
\begin{table}[h]
\small
\centering
\caption{SCIP heuristic results across problem families.}
\label{tab:scip_heuristics}
\begin{tabular}{llcccc}
\toprule
Family & Heuristic & Time (s) & Success (\%) & Rel.\ Subopt.\ (\%) & $\leq$5\%-optimal (\%) \\
\midrule
\multirow{3}{*}{CFLP-medium}    & FP & 7.4   & 100 & 2.64   & 93  \\
                                & RENS             & 3.8   & 100 & 0.09   & 100 \\
                                & SR  & 7.6   & 100 & 12.20  & 3   \\
\midrule
\multirow{3}{*}{GISP-easy}      & FP & 1.0   & 100 & 76.23  & 0   \\
                                & RENS             & 2.2   & 100 & 92.44  & 0   \\
                                & SR  & 1.0   & 100 & 94.92  & 0   \\
\midrule
\multirow{3}{*}{MIS-medium}     & FP & 11.0  & 100 & 122.91 & 0   \\
                                & RENS             & 10.6  & 89  & 100.01 & 0   \\
                                & SR  & 10.7  & 100 & 100.27 & 0   \\
\midrule
\multirow{3}{*}{MMCN-medium-BI} & FP & 2.1   & 0   & ---    & --- \\
                                & RENS             & 2.5   & 100 & 3.67   & 80  \\
                                & SR  & 2.1   & 2   & 135.63 & 0   \\
\midrule
\multirow{3}{*}{MVC-medium}     & FP & 19.1  & 10  & 37.38  & 0   \\
                                & RENS             & 18.5  & 10  & 89.72  & 0   \\
                                & SR  & 19.3  & 10  & 88.02  & 0   \\
\midrule
\multirow{3}{*}{OTS-medium}     & FP & 188.9 & 36  & 50.36  & 0   \\
                                & RENS             & 200.7 & 0   & ---    & --- \\
                                & SR  & 208.5 & 0   & ---    & --- \\
\bottomrule
\end{tabular}
\end{table}
}

\section{Tabular results} \label{sec:tables}
Time (s) refers to time until termination.
Suboptimality refers to the relative suboptimality of the returned solution.
Nodes refers to the number of nodes explored in the branch-and-bound tree.
Correct refers to the percentage of instances for which the method found an $\tol$-optimal solution within the time limit. If the 600\,s time limit is exceeded, the solver is allowed to return its current best solution.
Speedup (s) refers to the improvement in solve time compared to the baseline method --- either CP or GRB depending on the method.
Time, Suboptimality, and Nodes are reported as mean $\pm$ standard deviation.
    % {
    % \setlength{\tabcolsep}{3pt}
    % \centering
    % \small
% \textbf{EcoleIndset}
% \begin{tabular}{l l l l}
% \toprule
% Method & Time (s) & RelSubopt (\%) & Correct (\%) \\
% \midrule
% \csvreader[
%   head to column names,
%   late after line=\\
% ]{split_resultsCFLP.csv}{}%
% {%
%   \Method & \Time & \RelSubopt & \Pcorrect
% }
% \bottomrule
% \end{tabular}
\ifarxiv
\newpage
\fi
\vspace{1em}

\begin{table}[H]
\small
\centering
\caption{Capacitated facility location (CFLP-medium).}
\label{tab:results-cflp}
\begin{tabular}{llllll}
\toprule
Method & Time (s) & Suboptimality & Nodes & Correct & Speedup (s)\\
\midrule
\csvreader[
  head to column names=false,
  late after line=\\
]{CFLPaggregated_results_mean_pm_sd.csv}{}%
{\csvcoli & \csvcolii & \csvcoliv & \csvcolv & \csvcolvi & \csvcolvii}
\bottomrule
\end{tabular}
\end{table}

\begin{table}[H]
\small
\centering
\caption{Generalized independent set (GISP-easy).}
\label{tab:results-gisp}
\begin{tabular}{llllll}
\toprule
Method & Time (s) & Suboptimality & Nodes & Correct & Speedup (s)\\
\midrule
\csvreader[
  head to column names=false,
  late after line=\\
]{GISPaggregated_results_mean_pm_sd.csv}{}%
{\csvcoli & \csvcolii & \csvcoliv & \csvcolv & \csvcolvi & \csvcolvii}
\bottomrule
\end{tabular}
\end{table}
% \bnote{This ``Maximum vertex capacity'' block reads the MVC CSV and duplicates the ``Minimum vertex cover'' table further below. Decide whether to delete this copy or relabel/redirect to a different CSV.}
% \textbf{Maximum vertex capacity}
% \begin{table}[H]
% \small
% \centering
% \begin{tabular}{llllll}
% \toprule
% Method & Time (s) & Suboptimality & Nodes & Correct & Speedup (s)\\
% \midrule
% \csvreader[
%   head to column names=false,
%   late after line=\\
% ]{MVCaggregated_results_mean_pm_sd.csv}{}%
% {\csvcoli & \csvcolii & \csvcoliv & \csvcolv & \csvcolvi & \csvcolvii}
% \bottomrule
% \end{tabular}
% \end{table}

\ifarxiv
\newpage
\fi
\begin{table}[H]
\small
\centering
\caption{Optimal transmission switching (OTS-medium).}
\label{tab:results-ots}
\begin{tabular}{llllll}
\toprule
Method & Time (s) & Suboptimality & Nodes & Correct & Speedup (s)\\
\midrule
\csvreader[
  head to column names=false,
  late after line=\\
]{OTSaggregated_results_mean_pm_sd.csv}{}%
{\csvcoli & \csvcolii  & \csvcoliv & \csvcolv & \csvcolvi & \csvcolvii}
\bottomrule
\end{tabular}
\end{table}

\ifarxiv
\else
\newpage
\fi
\begin{table}[H]
\small
\centering
\caption{Middle-mile consolidation network design (MMCN-medium-BI).}
\label{tab:results-mmcn}
\begin{tabular}{llllll}
\toprule
Method & Time (s) & Suboptimality & Nodes & Correct & Speedup (s)\\
\midrule
\csvreader[
  head to column names=false,
  late after line=\\
]{MMCNaggregated_results_mean_pm_sd.csv}{}%
{\csvcoli & \csvcolii & \csvcoliv & \csvcolv & \csvcolvi & \csvcolvii}
\bottomrule
\end{tabular}
\end{table}

\begin{table}[H]
\small
\centering
\caption{Maximum independent set (MIS-medium).}
\label{tab:results-mis}
\begin{tabular}{llllll}
\toprule
Method & Time (s) & Suboptimality & Nodes & Correct & Speedup (s)\\
\midrule
\csvreader[
  head to column names=false,
  late after line=\\
]{MISaggregated_results_mean_pm_sd.csv}{}%
{\csvcoli & \csvcolii & \csvcoliv & \csvcolv & \csvcolvi & \csvcolvii}
\bottomrule
\end{tabular}
\end{table}

\begin{table}[H]
\small
\centering
\caption{Minimum vertex cover (MVC-medium).}
\label{tab:results-mvc}
\begin{tabular}{llllll}
\toprule
Method & Time (s) & Suboptimality & Nodes & Correct & Speedup (s)\\
\midrule
\csvreader[
  head to column names=false,
  late after line=\\
]{MVCaggregated_results_mean_pm_sd.csv}{}%
{\csvcoli & \csvcolii & \csvcoliv & \csvcolv & \csvcolvi & \csvcolvii}
\bottomrule
\end{tabular}
\end{table}

\change{
\section{Further simple baseline experiments} \label{appendix:moreexperiments}
To further demonstrate that our method learns nontrivial patterns in the convergence curves, we provide comparisons between our methods and some simple baseline models.
We consider two further simple models.
In the first (linear) the LSTM is replaced by a generalized linear model which is trained to predict $\log(\predgap_\param(t))$ from $\covars_\param(t)$, and the conformal prediction step is performed as normal.
In the second (CP UB rate) we replace the LSTM predictor with a heuristic which outputs the change in upper bound over the last 5 seconds. The conformal prediction is performed as normal after this.
We further compare our method to Gurobi with a larger range of tolerances, given by 1\%, 2\%, 5\%, and 10\% (0.1\% is given in the tables in Appendix~\ref{sec:tables}).
\subsection{Heuristic models}
Tabular results for the heuristic models are displayed in Table~\ref{tab:heur_subopt} (relative suboptimality) and Table~\ref{tab:heur_time} (solve time).
\begin{table}[h]
\small
\centering
\caption{Suboptimality (\%) by family for LSTM (ours) and the described heuristic models (mean $\pm$ std).}
\label{tab:heur_subopt}
\begin{tabular}{lcccc}
\toprule
Family & LSTM (ours) & Linear & CP UB rate & Full solve \\
\midrule
\csvreader[
  no head,
  late after line=\\,
  filter ifthen={\value{csvinputline}>1}
]{family_heuristic_subopt_table.csv}{}%
{\csvcoli & \csvcolii & \csvcoliii & \csvcoliv & \csvcolv}
\bottomrule
\end{tabular}
\end{table}
\begin{table}[h]
\small
\centering
\caption{Solve time by family for LSTM (ours) and the described heuristic models (mean $\pm$ std).}
\label{tab:heur_time}
\begin{tabular}{lcccc}
\toprule
Family & LSTM (ours) & Linear & CP UB rate & Full solve \\
\midrule
\csvreader[
  no head,
  late after line=\\,
  filter ifthen={\value{csvinputline}>1}
]{family_heuristic_time_table.csv}{}%
{\csvcoli & \csvcolii & \csvcoliii & \csvcoliv & \csvcolv}
\bottomrule
\end{tabular}
\end{table}
%
\subsection{Additional solver tolerances}
Results for a wider range of Gurobi solver tolerances (MIPGap) are displayed in Table~\ref{tab:family_subopt} (relative suboptimality) and Table~\ref{tab:family_time} (solve time).
Tolerances of 1\% and 2\% usually terminate after the optimal solution has been found, but on some problems they take significantly longer than our method and provide no statistical guarantees on optimality.
\begin{table}[h]
\footnotesize
\centering
\caption{Suboptimality (\%) by family for LSTM (ours) and Gurobi at various tolerances (mean $\pm$ std).}
\label{tab:family_subopt}
\begin{tabular}{lccccc}
\toprule
Family & LSTM (ours) & 1\% & 2\% & 5\% & 10\% \\
\midrule
\csvreader[
  no head,
  late after line=\\,
  filter ifthen={\value{csvinputline}>1}
]{family_subopt_table.csv}{}%
{\csvcoli & \csvcolii & \csvcoliii & \csvcoliv & \csvcolv & \csvcolvi}
\bottomrule
\end{tabular}
\end{table}
\begin{table}[h]
\footnotesize
\centering
\caption{Solve time by family for LSTM (ours) and Gurobi at various tolerances (mean $\pm$ std).}
\label{tab:family_time}
\begin{tabular}{lccccc}
\toprule
Family & LSTM (ours) & 1\% & 2\% & 5\% & 10\% \\
\midrule
\csvreader[
  no head,
  late after line=\\,
  filter ifthen={\value{csvinputline}>1}
]{family_time_table.csv}{}%
{\csvcoli & \csvcolii & \csvcoliii & \csvcoliv & \csvcolv & \csvcolvi}
\bottomrule
\end{tabular}
\end{table}
}

\change{
\section{Sensitivity analyses}
In this section we perform two additional experiments.
In the first we consider perturbations of training data to investigate how quickly performance degrades out of distribution.
In the second we investigate the effect of changing the size of the calibration set on coverage rate.
We only consider the parametric families CFLP, GISP, MIS and MVC in this section because the experiments require repeated resampling of iid datapoints.
This is not possible on the datasets which come from real data.
\subsection{Sensitivity to distribution shifts}
We used distributional MIPLIB generators for four families to investigate the effect of distribution shift.
On each problem which comes with a generator (CFLP, GISP, MIS and MVC), we generate a new train and calibration dataset, of size 1000 and 100 respectively, and generate 100 separate iid test sets of size 100.
In each experiment we perturb all random continuous parameters in the generator of the test dataset by multiplying by a uniform random variable $U[1, 1+p]$, where $p$ controls the perturbation magnitude.
We perform the conformal procedure as described in the main text.
Results are displayed in Table~\ref{tab:perturbation_subopt}.
We observe no systematic degradation in suboptimality, consistent with the TV-distance guarantee.

\begin{table}[h]
\small
\centering
\caption{Suboptimality (\%) under distribution shift (mean $\pm$ std).}
\label{tab:perturbation_subopt}
\begin{tabular}{lcccc}
\toprule
Family & $p=0$ & $p=0.01$ & $p=0.05$ & $p=0.2$ \\
\midrule
CFLP & $0.020 \pm 0.059$ & $0.023 \pm 0.046$ & $0.009 \pm 0.027$ & $0.010 \pm 0.031$ \\
GISP & $0.055 \pm 0.462$ & $0.135 \pm 0.749$ & $0.091 \pm 0.574$ & $0.075 \pm 0.327$ \\
MIS  & $0.007 \pm 0.032$ & $0.006 \pm 0.026$ & $0.012 \pm 0.064$ & $0.004 \pm 0.014$ \\
MVC  & $0.031 \pm 0.071$ & $0.016 \pm 0.043$ & $0.013 \pm 0.033$ & $0.024 \pm 0.051$ \\
\bottomrule
\end{tabular}
\end{table}

\subsection{Sensitivity to calibration set size}
On each problem which comes with a generator (CFLP, GISP, MIS and MVC), we generate a new train and test dataset, of size 1000 and 100 respectively, and generate 100 separate iid calibration sets of size 100.
The conformal procedure is carried out as described in the main paper, and results are reported on the test dataset.
Results are displayed in Table~\ref{tab:calib_size}.
Note that because we have regenerated the problems ourselves, to allow repeated resampling of the data, the solve-times do not always line up with the solve-times of problems from the Distributional MIPLIB dataset. 
For some problems, we were unable to figure out exactly which parameters were used to generate the data.
However, the results tell the expected story. Variance of coverage is very high when small calibration sets are used, and it decreases as the size of the calibration set increases.
Solve time increases on average as the calibration set size increases, but coverage is more reliable.
\begin{table}[ht]
\centering
\caption{Effect of the calibration set size $c$ on coverage and solve time.}
\label{tab:calib_size}
\ifarxiv
\begin{subtable}{\textwidth}
\else
\begin{subtable}{0.48\textwidth}
\fi
\centering
\small
\setlength{\tabcolsep}{4pt}
\caption{Coverage rate (\%).}
\label{tab:calib_coverage}
\begin{tabular}{lcccc}
\toprule
Family & $c=5$ & $c=10$ & $c=20$ & $c=100$ \\
\midrule
CFLP & $80.0 \pm 16.4$ & $87.9 \pm 9.1$  & $88.8 \pm 8.6$ & $92.7 \pm 2.7$ \\
GISP & $75.1 \pm 18.2$ & $80.5 \pm 12.0$ & $89.7 \pm 7.4$ & $92.8 \pm 3.3$ \\
MIS  & $77.5 \pm 12.5$ & $86.8 \pm 5.6$  & $91.4 \pm 5.5$ & $95.3 \pm 3.0$ \\
MVC  & $77.7 \pm 15.4$ & $85.3 \pm 11.3$ & $90.5 \pm 7.5$ & $96.0 \pm 1.5$ \\
\bottomrule
\end{tabular}
\end{subtable}\ifarxiv\\\medskip\else\hfill\fi
\ifarxiv
\begin{subtable}{\textwidth}
\else
\begin{subtable}{0.48\textwidth}
\fi
\centering
\small
\setlength{\tabcolsep}{4pt}
\caption{Mean solve time (s).}
\label{tab:calib_time}
\begin{tabular}{lcccc}
\toprule
Family & $c=5$ & $c=10$ & $c=20$ & $c=100$ \\
\midrule
CFLP & $9.4 \pm 4.6$  & $11.3 \pm 3.2$ & $11.4 \pm 3.0$  & $13.1 \pm 1.6$ \\
GISP & $20.7 \pm 7.0$ & $22.0 \pm 5.6$ & $26.8 \pm 3.9$  & $28.1 \pm 2.0$ \\
MIS  & $8.7 \pm 1.1$  & $10.1 \pm 2.3$ & $15.1 \pm 13.1$ & $14.7 \pm 4.3$ \\
MVC  & $8.8 \pm 0.9$  & $9.3 \pm 0.6$  & $9.6 \pm 0.3$   & $9.7 \pm 0.0$  \\
\bottomrule
\end{tabular}
\end{subtable}
\end{table}
}
% \textbf{Optimal transmission scheduling}\\
% \begin{tabular}{l l l l}
% \toprule
% Method & Time (s) & RelSubopt (\%) & Correct (\%) \\
% \midrule
% \csvreader[
%   head to column names,
%   late after line=\\
% ]{split_resultsOTS.csv}{}%
% {%
%   \Method & \Time & \RelSubopt & \Pcorrect
% }
% \bottomrule
% \end{tabular}

% \textbf{Middle-mile consolidation network design}\\
% \begin{tabular}{l l l l}
% \toprule
% Method & Time (s) & RelSubopt (\%) & Correct (\%) \\
% \midrule
% \csvreader[
%   head to column names,
%   late after line=\\
% ]{split_resultsMMCN.csv}{}%
% {%
%   \Method & \Time & \RelSubopt & \Pcorrect
% }
% \bottomrule
% \end{tabular}

% \textbf{Maximum independent set}\\
% \begin{tabular}{l l l l}
% \toprule
% Method & Time (s) & RelSubopt (\%) & Correct (\%) \\
% \midrule
% \csvreader[
%   head to column names,
%   late after line=\\
% ]{split_resultsMIS.csv}{}%
% {%
%   \Method & \Time & \RelSubopt & \Pcorrect
% }
% \bottomrule
% \end{tabular}

% \textbf{Maximum vertex cover}\\
% \begin{tabular}{l l l l}
% \toprule
% Method & Time (s) & RelSubopt (\%) & Correct (\%) \\
% \midrule
% \csvreader[
%   head to column names,
%   late after line=\\
% ]{split_resultsMVC.csv}{}%
% {%
%   \Method & \Time & \RelSubopt & \Pcorrect
% }
% \bottomrule
% \end{tabular}

\end{document}




% This document was modified from the file originally made available by
% Pat Langley and Andrea Danyluk for ICML-2K. This version was created
% by Iain Murray in 2018, and modified by Alexandre Bouchard in
% 2019 and 2021 and by Csaba Szepesvari, Gang Niu and Sivan Sabato in 2022.
% Modified again in 2023 and 2024 by Sivan Sabato and Jonathan Scarlett.
% Previous contributors include Dan Roy, Lise Getoor and Tobias
% Scheffer, which was slightly modified from the 2010 version by
% Thorsten Joachims & Johannes Fuernkranz, slightly modified from the
% 2009 version by Kiri Wagstaff and Sam Roweis's 2008 version, which is
% slightly modified from Prasad Tadepalli's 2007 version which is a
% lightly changed version of the previous year's version by Andrew
% Moore, which was in turn edited from those of Kristian Kersting and
% Codrina Lauth. Alex Smola contributed to the algorithmic style files.
